%%%%%%%%%%%%%%%%%%%%%%%%%%%%%%%%%%%%%%%%%%%%%%%%%%%%%%%%%%%%%%%%%%%%%%
% Câu chưa có mã dạng — cần gắn %[..] rồi chuyển vào đúng file dạng
% Số câu: 36  (sinh tự động bởi tach_cau.py; từ nay thầy sửa trực tiếp file này)
%%%%%%%%%%%%%%%%%%%%%%%%%%%%%%%%%%%%%%%%%%%%%%%%%%%%%%%%%%%%%%%%%%%%%%

%%%%%%%%%%%%%%%% TRẮC NGHIỆM 4 PHƯƠNG ÁN %%%%%%%%%%%%%%%%

% nguồn: Phan1_Lop_12_Hinh_hoc_Chuong_5
\begin{ex}%[ma:4e8847df]
	]%[2H5H1-6] 
	 Trong không gian $Oxyz$, cosin của góc giữa hai mặt phẳng $(P)\colon x + 2y - 2z + 1 = 0$ và $(Q)\colon x + y + z - 1 = 0$ bằng
	\choice 
		 {$\dfrac{\sqrt{3}}{3}$}
		{$-\dfrac{\sqrt{3}}{3}$} 
		{$-\dfrac{\sqrt{3}}{9}$} 
		{\True $\dfrac{\sqrt{3}}{9}$} 
	\loigiai{Hai  mặt phẳng $(P)$ và $(Q)$ có các vectơ pháp tuyến lần lượt là $\overrightarrow{n}_P = (1;2;-2)$ và $\overrightarrow{n}_Q = (1;1;1)$.\\
	Gọi $\phi$ là số đo của góc  giữa $(P)$ và $(Q)$, ta có\\
	$\cos \phi = \dfrac{\big|\overrightarrow{n}_P \cdot \overrightarrow{n}_Q\big|}{\big|\overrightarrow{n}_P\big| \cdot \big|\overrightarrow{n}_Q\big|} = \dfrac{\big|1\cdot1 + 2\cdot1 + (-2)\cdot1\big|}{\sqrt{1^2+2^2+(-2)^2} \cdot \sqrt{1^2+1^2+1^2}} =  \dfrac{1}{3\cdot\sqrt{3}} = \dfrac{\sqrt{3}}{9}$.
	} 
\end{ex}

% nguồn: Phan1_Lop_12_Hinh_hoc_Chuong_5
\begin{ex}%[ma:03d7b654]
	%[2H5N1-2] 
	Cho hình hộp chữ nhật $ABCD.A'B'C'D'$. Cặp vectơ nào sau đây~là cặp vectơ chỉ phương của mặt phẳng  $\left(ABB'A'\right)$ ?
	 \choice
		{$\overrightarrow{AB}$  và $\overrightarrow{AD}$}
		 {$\overrightarrow{AB}$ và $\overrightarrow{AD'}$}
		{$\overrightarrow{AB}$ và  $\overrightarrow{A'B'}$}
		{\True $\overrightarrow{AB}$ và  $\overrightarrow{CC'}$}
	\loigiai{ 
	\immini{ 
		Vì hai vectơ $\overrightarrow{AB}$ và  $\overrightarrow{C C'}$ thỏa mãn được các điều kiện sau\\
		 a) Là các vectơ khác vectơ $\overrightarrow{0}$. \\
		b) Không cùng phương.  \\
		c) Vectơ $\overrightarrow{A B}$ có giá~là đường thẳng $AB$ nằm trong mặt phẳng $\left(ABB'A'\right)$ và  vectơ $\overrightarrow{C C'}$ có giá~là đường thẳng $CC'$ song song với mặt phẳng $\left(ABB'A'\right)$.\\
		Nên cặp vectơ $\overrightarrow{A B}$ và $\overrightarrow{C C'}$~là cặp vectơ chỉ  phương của mặt phẳng $\left(ABB'A'\right)$.
		}{\begin{tikzpicture}[line cap=round,line join=round]
		 \coordinate (A') at (1,1);
		\coordinate (B') at (0,0); 
		\coordinate  (C') at (3,0);
		\coordinate (D') at ($(A')-(B')+(C')$); 
		\coordinate (A) at (1,3); 
		\coordinate (B)  at ($(A)-(A')+(B')$);
		\coordinate (C)  at ($(A)-(A')+(C')$);
		\coordinate (D)  at ($(A)-(A')+(D')$);
		 \foreach \x/\g in{A'/45,B'/-90,C'/-90,D'/0,A/90,B/90,C/90,D/90}
		\fill (\x)  circle (1pt) node[shift={(\g:10pt)}]{$\x$};
		\draw (B')--(C') 
			(C')--(D') 
			(A)--(B)--(C)--(D)--(A) 
			 (B')--(B)
			(C')--(C) 
			(D')--(D); 
		 \draw[dashed] (A')--(A)
			 (A')--(B')
			 (A')--(D');
	\end{tikzpicture}
		} 
	} 
\end{ex}

% nguồn: Phan1_Lop_12_Hinh_hoc_Chuong_5
\begin{ex}%[ma:83ab5adf]
	 %[2H5N3-1]
	Trong không gian $Oxyz$, điểm  nào trong các điểm sau đây thuộc mặt cầu
	\[(S)\colon x^2  + y^2 + z^2=50?\]
	 \choice
		{$M(3; 4; 6)$} 
		{$N(4; 4; 5)$} 
		{\True $P(3;  4; -5)$}
		{$Q(- 3;  3; -5)$}
	 \loigiai{
		Thay toạ độ các điểm $M$, $N$, $P$, $Q$vào  phương trình mặt cầu. \\
		Thay $\heva{&x=3\\&y=4\\&z=6}$  vào phương trình mặt cầu, ta được $3^2+4^2+6^2=50$ (sai).\\
		Thay $\heva{&x=4\\&y=4\\&z=5}$ vào phương trình mặt cầu, ta được $4^2+4^2+5^2=50$  (sai).\\
		Thay $\heva{&x=3\\&y=4\\&z=-5}$ vào phương trình mặt cầu,  ta được $3^2+4^2+(-5)^2=50$ (đúng).\\
		Thay $\heva{&x=-3\\&y=3\\&z=-5}$ vào  phương trình mặt cầu, ta được $(-3)^2+3^2+(-5)^2=50$ (sai).\\
		 Chỉ có toạ độ điểm $P$ thỏa mãn.
	 }
\end{ex}

% nguồn: Phan1_Lop_12_Hinh_hoc_Chuong_5
\begin{ex}%[ma:d3c67d15]
	%[2H5H1-3] 
	Trong không gian $Oxyz$, mặt phẳng đi qua điểm $K(1; 1; 1)$ nhận $\vec{u}=(1; 0; 1)$, $\vec{v}=(1; 1; 0)$~là cặp vectơ chỉ phương có phương trình tổng quát~là 
	 \choice
		{$x+y+z-3=0$} 
		 {$x-y+z-1=0$}
		{$x+y-z-1=0$} 
		{\True $-x+y+z-1=0$} 
	\loigiai{ 
		Gọi $(P)$~là mặt phẳng đi qua điểm $K(1; 1; 1)$ nhận $\vec{u}=(1; 0; 1)$, $\vec{v}=(1; 1; 0)$~là cặp vectơ chỉ phương có dạng phương trình tổng quát $  ax+by+cz+d=0$.\\
		Vì $(P)$ nhận  cặp vectơ $\vec{u}$ và $\vec{v}$~là cặp vectơ chỉ phương nên ta có thể chọn được cho $(P)$ một vectơ pháp tuyến~là $\left[\vec{u}, \vec{v}\right]= (-1; 1; 1)$.\\
		Do đó $a=-1$, $b=1$, $c=1$. Mặt khác $(P)$ đi  qua điểm $K(1; 1; 1)$ suy ra $-1\cdot 1+1 \cdot 1+1 \cdot 1+d=0$ suy ra $d=-1$.\\
		Vậy $(P)$  có phương trình~là $-x+y+z-1=0$.
	 }
\end{ex}

% nguồn: Phan1_Lop_12_Hinh_hoc_Chuong_5
\begin{ex}%[ma:64822e0c]
	 %[2H5B1-2]
	Trong không gian với hệ tọa~độ $Oxyz$, góc $\varphi$ giữa hai mặt phẳng nhận $\vec{n}=\left(a;b;c \right)$,  $\vec{n}'=\left(a';b';c' \right)$~là vectơ pháp tuyến thỏa mãn
	 \choice
		{$\cos  \varphi =\dfrac{aa'+bb'+cc'}{\sqrt{a^2+b^2+c^2}\cdot\sqrt{{{{a'}}^2}+{{{b'}}^2}+{{{c'}}^2}}}$}
		{$\sin \varphi =\dfrac{aa'+bb'+cc'}{\sqrt{a^2+b^2+c^2}\cdot\sqrt{{{{a'}}^2}+{{{b'}}^2}+{{{c'}}^2}}}$} 
		{$\sin \varphi =\dfrac{\left|  aa'+bb'+cc' \right|}{\sqrt{a^2+b^2+c^2}\cdot\sqrt{{{{a'}}^2}+{{{b'}}^2}+{{{c'}}^2}}}$}
		{\True $\cos \varphi =\dfrac{\left| aa'+bb'+cc' \right|}{\sqrt{a^2+b^2+c^2}\cdot\sqrt{{{{a'}}^2}+{{{b'}}^2}+{{{c'}}^2}}}$} 
	 \loigiai{
		Góc $\varphi$ giữa hai mặt phẳng nhận $\vec{n}=\left(a;b;c \right)$,  $\vec{n}'=\left(a';b';c' \right)$~là vectơ pháp tuyến thỏa mãn biểu thức: $$\cos \varphi =\dfrac{\left| aa'+bb'+cc' \right|}{\sqrt{a^2+b^2+c^2}\cdot\sqrt{{{{a'}}^2}+{{{b'}}^2}+{{{c'}}^2}}}.$$
	} 
\end{ex}

% nguồn: Phan1_Lop_12_Hinh_hoc_Chuong_5
\begin{ex}%[ma:bbee851d]
	 %[2H5H2-2]
	Trong không  gian với hệ tọa~độ $ Oxyz$, góc $\alpha $ giữa đường thẳng nhận ${\vec u}\left(a;b;c\right)$~là véc-tơ chỉ phương và mặt phẳng nhận $\overrightarrow{n'}\left(a';b';c'\right)$~là véc-tơ pháp tuyến thoả mãn
	 \choice
		 {$\cos \alpha=\dfrac{\left|aa'+bb'+cc'\right|}{\sqrt{a^2+b^2+c^2}\cdot\sqrt{a'^2+b'^2+c'^2}}$}
		 {$\sin\alpha=\dfrac{aa'+bb'+cc'}{\sqrt{a^2+b^2+c^2}\cdot\sqrt{a'^2+b'^2+c'^2}}$}
		{\True $\sin\varphi=\dfrac{\left|aa'+bb'+cc'\right|}{\sqrt{a^2+b^2+c^2}\cdot\sqrt{a'^2+b'^2+c'^2}}$} 
		{$\cos  \varphi=\dfrac{aa'+bb'+cc'}{\sqrt{a^2+b^2+c^2}\cdot\sqrt{a'^2+b'^2+c'^2}}$}
	 \loigiai{
		Góc $\alpha $ giữa đường thẳng nhận ${\vec u}\left(a;b;c\right)$~là véc-tơ chỉ phương và mặt phẳng nhận $\overrightarrow{n'}\left(a';b';c'\right)$~là  véc-tơ pháp tuyến thoả mãn $\sin\varphi=\dfrac{\left|aa'+bb'+cc'\right|}{\sqrt{a^2+b^2+c^2}\cdot\sqrt{a'^2+b'^2+c'^2}}$.
	 }
\end{ex}

% nguồn: Phan1_Lop_12_Hinh_hoc_Chuong_5
\begin{ex}%[ma:3a2b9294]
	%[2H5H2-7] 
	Trong không gian với hệ tọa~độ $ Oxyz$, góc $\varphi $  giữa hai đường thẳng nhận ${\vec u}\left(a;b;c\right), \vec{u'}\left(a';b';c'\right)$~là véc-tơ chỉ phương thỏa mãn
	 \choice
		 {\True $\cos\varphi=\dfrac{\left|aa'+bb'+cc'\right|}{\sqrt{a^2+b^2+c^2}\cdot\sqrt{a'^2+b'^2+c'^2}}$}
		{$\sin\varphi=\dfrac{aa'+bb'+cc'}{\sqrt{{\sqrt{a^2+b^2+c^2}\cdot\sqrt{a'^2+b'^2+c'^2}}}}$} 
		{$\sin\varphi=\dfrac{\left|aa'+bb'+cc'\right|}{\sqrt{a^2+b^2+c^2}\cdot\sqrt{a'^2+b'^2+c'^2}}$} 
		{$\cos \varphi=\dfrac{aa'+bb'+cc'}{\sqrt{a^2+b^2+c^2}\cdot\sqrt{a'^2+b'^2+c'^2}}$} 
	\loigiai{Góc $\varphi $ giữa hai đường thẳng  nhận ${\vec u}\left(a;b;c\right), \vec{u'}\left(a';b';c'\right)$~là véc-tơ chỉ phương thỏa mãn
		 $$\cos\varphi=\dfrac{\left|aa'+bb'+cc'\right|}{\sqrt{a^2+b^2+c^2}\cdot\sqrt{a'^2+b'^2+c'^2}}$$
	} 
\end{ex}

% nguồn: Phan1_Lop_12_Hinh_hoc_Chuong_5
\begin{ex}%[ma:0aabf9f1]
	 %[2H5H2-3]
	Trong không  gian $Oxyz$, đường thẳng đi qua điểm $I(15;-16; 17)$ và nhận $\vec{u}=(-7; 8;-9)$~là vectơ chỉ phương có phương trình tham số~là
	 \choice
		{\True $\heva{&x=15-7  t \\& y=-16+8 t \\& z=17-9 t}$}
		{$\heva{&x=15-7  t \\& y=-16+8 t \\& z=17-9 t^2}$}
		{$\heva{&x=15-7 t^2  \\& y=-16+8 t \\& z=17-9 t}$}
		{$\heva{&x=-7+15 t \\& y=8-16 t \\& z=-9+17 t}$} 
	 \loigiai{
		Đường thẳng đi qua điểm $I(15;-16; 17)$ và  nhận $\vec{u}=(-7; 8;-9)$~làm vectơ chỉ phương có phương trình tham số~là $\heva{&x=15-7 t \\& y=-16+8 t \\& z=17-9 t.}$
	} 
\end{ex}

% nguồn: Phan1_Lop_12_Hinh_hoc_Chuong_5
\begin{ex}%[ma:6bc695dc]
	%[2H5H2-4] 
	Cho hai đường thẳng $d\colon  \heva{& x=2+t\\ &y=1-t\\ &z=1}$ $(t \in \mathbb{R})$ và $d'\colon \dfrac{x-1}{1}=\dfrac{y-2}{-1}=\dfrac{z+2}{-2}$. Khẳng định nào sau đây đúng?
	 \choice
		{$d$, $d'$ song song với nhau} 
		{$d$, $d'$ chéo nhau} 
		{$d$, $d'$ trùng nhau} 
		{\True $d$, $d'$ cắt nhau} 
	\loigiai{ 
		Hai đường thẳng $d$ và $d'$ có véc-tơ chỉ phương lần lượt~là $\vec{a}=(1;-1;0)$ và \break $\vec{a}'=(1;-1;-2)$.\\ 
		Ta  có $\vec{a}$ và $\vec{a}'$ không cùng phương nên $d$ và $d'$ hoặc cắt nhau hoặc chéo nhau.\\
		Đường thẳng $d'$ có phương trình  tham số~là $d'\colon \heva{& x=1+t'\\ &y=2-t'\\ &z=-2-2t'}\quad\left(t' \in \mathbb{R}\right)$.\\
		 Giải hệ phương trình $\heva{& 2+t=1+t'\\ &1-t=2-t'\\ &1=-2-2t'}$ ta được $t=-\dfrac{5}{2}$, $t'=-\dfrac{3}{2}$.\\
		Vậy $d$ cắt $d'$. 
	} 
\end{ex}


%%%%%%%%%%%%%%%% ĐÚNG/SAI %%%%%%%%%%%%%%%%

% nguồn: Phan1_Lop_12_Hinh_hoc_Chuong_5
\begin{ex}%[ma:6490e64d]
	 %[2H5H2-8]
	Hệ thống Kiểm soát không lưu, còn gọi là~kiểm soát không lưu (tiếng anh: air traffic control, viết tắt là~ATC), hay Điều khiển không lưu là~hệ thống chuyên trách đảm nhận việc gửi các hướng dẫn đến máy bay nhằm giúp các máy bay tránh va chạm, đồng thời đảm bảo tính hoạt động hiệu quả của nền tảng không lưu. Trong không gian với hệ tọa độ $Oxyz$, xét một đài kiểm soát không lưu sân bay có tọa độ $O(0;0;0)$, mỗi đơn vị trên trục ứng với 1 km. Máy bay bay trong phạm vi cách đài kiểm soát 417 km sẽ hiển thị trên màn hình ra đa. Một máy bay đang ở  vị trí $A(-688; -185; 8)$ chuyển động theo đường thẳng $d$ có vectơ chỉ phương là~$\overrightarrow{u} = (91; 75; 0)$ và hướng về đài kiểm soát không lưu. Xét các phát biểu sau:
	\begin{center}
		\begin{tikzpicture}[line join=round,line cap=round,scale=0.7]
				 \path
					 (-6,2.1) coordinate (M)
					(8,2.1) coordinate  (N)
					(-8,-3.5)  coordinate (Q)
					($(N)+(Q)-(M)$) coordinate (P) 
					 (0,0) coordinate (O)
					($(O)-(2.5,0)$) coordinate (H) 
					(-5,1.5) coordinate (A) 
					(3.5,1) coordinate  (D);
				\path[name path=d] (A)--(D) node[above]{$  d $};
				\path[name path=(C)] (O) circle (2.5  cm);
				\path[name path=(E)] (O) ellipse (2.5 cm and 1.5  cm);
				 \path[name intersections={of=d and (E),by={C,B}}];
				\path[name intersections={of=d and (C),by={I,J}}]; 
				 \draw (A)--(J)
					(I)--(D); 
				\draw[dashed]  (J)--(I);
				 \draw
					 (H) arc (180:0:2.5)
					(H) arc (-180:0:2.5 cm and 1.5 cm); 
				 \draw[dashed,thin]
					(H) arc (180:0:2.5 cm  and 1.5 cm)
					(H)  arc (-180:0:2.5);
				\path[name path=MN]  (M)--(N);
				\path[name intersections={of=MN  and (C),by={E,F}}];
				\draw[dashed]  (E)--(F);
				\draw (M)--(F) 
					(E)--(N)--(P)--(Q)--(M); 
				\draw[dashed,->]  ($(O)+(2*2/3,3*2/3)$)--($(O)-(2,3)$) node[left]{$ x $}; \draw (-0.94,-1.4)--(-2,-3);
				\draw[dashed,->] ($(O)-(4,0)$)--($(O)+(4,0)$) node[below]{$  y $}; \draw (-4,0)--(-2.5,0)
					(2.5,0)--(4,0); 
				\draw[dashed,->] ($(O)-(0,3)$)--($(O)+(0,3.5)$) node[left]{$  z $}; \draw (0,2)--(0,3.5);
				\foreach \i/\g in {O/-40,B/-90,C/-90,A/90}{\fill (\i)  circle (1pt)
					 ($(\i)+(\g:3.3mm)$) node[scale=1]{$\i$};}
			\end{tikzpicture}
	\end{center}
	\choiceTF 
		{\True Phương trình đường thẳng mô tả đường bay của máy bay trên là: $ \heva{&x = -688 - 91t\\ &y  = -185 + 75t\\ &z = 8},\ t\in \mathbb{R} $}
		{\True Xác định tọa độ của vị trí mà máy bay bay gần  đài kiểm soát không lưu nhất là~điểm $\left(-\dfrac{375}{2}, \dfrac{455}{2}, 8\right)$}
		{Vị trí sớm nhất mà máy bay  xuất hiện trên màn hình ra đa có tọa độ $(-88; 415; 8)$}
		{Giả sử suốt quá trình được theo dõi bởi đài kiểm soát không lưu này  máy bay luôn giữ vận tốc không đổi là~800 km/h thì mất 0,62 giờ (làm tròn đến hàng phần trăm)}
	\loigiai{ 
		\begin{itemchoice}
			\itemch Phương trình tham số của đường thẳng $d$  đi qua điểm $A(-688; -185; 8)$ và có vectơ chỉ phương $\overrightarrow{u} = (91; 75; 0)$ là:
			$ \heva{&x = -688 + 91t  \\ &y = -185 + 75t \\ &z = 8}\ (t \text{ là~tham số}) $.
			\itemch Gọi $H$ là~vị trí mà máy bay bay  gần đài kiểm soát không lưu nhất. Khi đó, khoảng cách $OH$ phải ngắn nhất, điều này xảy ra khi và chỉ khi $OH \perp d$.\\
			Vì $H  \in d$ nên $H(-688+91t;-185+75t;8)$.\\
			Ta có $\overrightarrow{OH} =  (-688+91t;-185+75t;8)$.\\
			$OH \perp d \Leftrightarrow \overrightarrow{OH}.\overrightarrow{u} = 0  \Leftrightarrow (-688+91t).91 + (-185+75t).75 + 8.0 = 0$\\
			$\Leftrightarrow 13906t - 76483 = 0 \Leftrightarrow t = \dfrac{11}{2}$\\ 
			Do đó  $H \left(-\dfrac{375}{2};\dfrac{455}{2};8\right)$.
			\itemch Lưu ý rằng, $M$ bay từ $A$ hướng về $ATC$ khi và chỉ khi $t$  lớn dần từ 0.\\
			Gọi $B$ là~vị trí sớm nhất mà máy bay xuất hiện trên màn hình ra  đa.\\
			Vì $B  \in d$ nên $B(-688+91t;-185+75t;8)$.\\
			$B$ là~vị trí sớm nhất mà máy bay xuất hiện  trên màn hình ra đa khi $OB = 417$.\\
			 $\Leftrightarrow \sqrt{(-688+91t)^2 + (-185+75t)^2 + 8^2} = 417$.\\
			$\Leftrightarrow 13906t^2 - 152966t + 333744 = 0$.\\ 
			$\Leftrightarrow t =  3$ hoặc $t = 8$ (loại $t=8$).\\
			Với $t  = 3$, ta có $B(-415; 40; 8)$.
			\itemch Vị trí đầu tiên và cuối cùng của $M$ trên màn hình ra đa là~$B(-415; 40; 8)$ và $C(40; 415; 8)$. Do đó thời  gian xuất hiện trên màn hình ra đa là
			$$\frac{BC}{800} = \frac{\sqrt{(40+415)^2 + (415-40)^2 + (8-8)^2}}{800}  \approx 0,74 \text{ (giờ)}.$$
		\end{itemchoice}
	} 
\end{ex}

% nguồn: Phan1_Lop_12_Hinh_hoc_Chuong_5
\begin{ex}%[ma:1c5951a6]
	 %[2H5H2-7]
	Trong không gian  với hệ tọa~độ $Oxyz$, cho đường thẳng $\Delta\colon \dfrac{x-2}{2}=\dfrac{y+8}{3}=\dfrac{z-7}{6}$ và mặt phẳng $(P)\colon x + 2y+ 2z + 2025 = 0$.
	\choiceTF 
		{$\vec{u}=(2; 6;  3)$~là một vectơ chỉ phương của $\Delta$}
		{\True$\vec{n}=(1; 2;  2)$~là một vectơ pháp tuyến của $(P)$}
		{\True $\vec{u}\cdot \vec{n}=20$ với $\vec{u}=(2; 3; 6)$~là một vectơ  chỉ phương của $\Delta$}
		{\True Góc giữa đường  thẳng $\Delta$ và mặt phẳng $(P)$ (làm tròn đến hàng đơn~vị của~độ)~là $72^\circ$}
	 \loigiai{
		\begin{itemchoice}
			\itemch \textbf{Sai.} Vì đường thẳng $\Delta\colon \dfrac{x-2}{2}=\dfrac{y+8}{3}=\dfrac{z-7}{6}$ có một  vectơ chỉ phương $\vec{u}_{\Delta}=(2; 3; 6)$.
			 \itemch \textbf{Đúng.} Vì mặt phẳng $(P)\colon x + 2y+ 2z + 2025 = 0$ có một vectơ pháp tuyến $\vec{n}=(1; 2; 2)$.
			\itemch \textbf{Đúng.}  Vì ta có $\vec{u}_{\Delta}\cdot \vec{n}=1\cdot2+2\cdot3+2\cdot6=20$.
			\itemch \textbf{Đúng.} Vì  $\sin\left((P),\Delta\right) =\dfrac{\left| \vec{n}\cdot\vec{u}_{\Delta}\right| }{\left|\vec{n}\right|\cdot\left|\vec{u}_{\Delta}\right|}=\dfrac{\left|20\right|}{\sqrt{49}\cdot \sqrt{9}}=\dfrac{20}{21}\Rightarrow \left((P),\Delta\right)\approx 72^\circ$.
		\end{itemchoice}
	 }
\end{ex}

% nguồn: Phan1_Lop_12_Hinh_hoc_Chuong_5
\begin{ex}%[ma:df1ff70c]
	 %[2H5N2-2]%[2H2H2-5]%[2H5H2-4]
	Trong  không gian $Oxyz$ cho hai đường thẳng
	\begin{align*}
		 \Delta_1\colon \dfrac{x}{1}=\dfrac{y-3}{-1}=\dfrac{z+3}{2},
		 \Delta_2\colon \dfrac{x+4}{2}=\dfrac{y+2}{1}=\dfrac{z-4}{-1}.
	\end{align*}
	Xét  các vectơ $\vec{u}_1=(1;-1;2)$ và $\vec{u}_2=(2;1;-1)$.
	\choiceTF 
		{\True  Đường thẳng $\Delta_1$ đi qua điểm $M_1(0;3;-3)$ và có $\vec{u}_1=(1;-1;2)$~làm một vectơ chỉ phương}
		{\True Đường thẳng $\Delta_2$ đi qua  điểm $M_2(-4;-2;4)$ và có $\vec{u}_2=(2;1;-1)$~làm một vectơ chỉ phương}
		 {$\left[\vec{u}_1,\vec{u}_2\right]=(1;-5;3)$}
		{Hai đường thẳng $\Delta_1$ và $\Delta_2$ chéo nhau} 
	\loigiai{ 
		\begin{itemchoice}
			\itemch  \textbf{Đúng}.\\
			Đường thẳng $\Delta_1\colon \dfrac{x}{1}=\dfrac{y-3}{-1}=\dfrac{z+3}{2}$ đi qua điểm $M_1(0;3;-3)$ có một vectơ chỉ  phương~là $\vec{u}_1=(1;-1;2)$.
			\itemch  \textbf{Đúng}.\\
			Đường thẳng $\Delta_2\colon \dfrac{x+4}{2}=\dfrac{y+2}{1}=\dfrac{z-4}{-1}$ đi qua điểm $M_2(-4;-2;4)$ có một  vectơ chỉ phương~là $\vec{u}_2=(2;1;-1)$.
			\itemch  \textbf{Sai}.\\
			Ta có $\heva{\vec{u}_1&=(1;-1;2)\\ \vec{u}_2&=(2;1;-1)}$\\ 
			Tích có hướng $\left[\vec{u}_1,\vec{u}_2\right] 
			=\left(\begin{vmatrix}
			-1 &  2 \\ 1 & -1
		\end{vmatrix};
		\begin{vmatrix}
			2  & 1 \\ -1 & 2
		\end{vmatrix};
		\begin{vmatrix}
			1 & -1 \\ 2  & 1
		\end{vmatrix}\right)=(-1;5;3)$.
		 \itemch \textbf{Sai}.\\
		 Ta có $\heva{&\vec{u}_1=(1;-1;2)\\& \vec{u}_2=(2;1;-1)\\& \vec{M_1M}_2=(-4;-5;7)}$\\
		 Xét tích vô hướng $\left[\vec{u}_1,\vec{u}_2\right]\cdot\vec{M_1M}_2 =-1\cdot(-4)+5\cdot(-5)+3\cdot7=0$.\\
		Vậy hai đường thẳng $\Delta_1$ và $\Delta_2$ không chéo  nhau.
	\end{itemchoice}
} 
\end{ex}

% nguồn: Phan1_Lop_12_Hinh_hoc_Chuong_5
\begin{ex}%[ma:35376085]
	%[2H5V1-7] 
	Các thiên thạch có đường kính lớn hơn $140$ m và có thể lại gần Trái Đất ở khoảng cách nhở hơn $7\,500\,000$ km được coi là~những vật thể có khả năng và chạm gây nguy hiểm cho Trái Đất. Để theo dõi những thiên thạch này, người ta đã thiết lập các trạm quan sát các vật thể bay gần Trái Đất. Giả sử có một hệ thống quan sát có khả năng theo dõi các vật thể ở độ cao không vượt quá $6630$ km so với mực nước biển. Coi Trái Đất là~khối cầu có bán  kính $6370$ km. Chọn hệ trục toạ độ $Oxyz$ trong không gian có gốc $O$ tại tâm Trái Đất và đơn vị độ dài trên mỗi trục tọa độ là~$1000$ km. Một thiên thạch (coi như một hạt) chuyển động với tốc độ không đổi theo một đường thẳng từ điểm $M\left(6; 15; -2 \right)$ sau một thời gian vị trí đầu tiên thiên thạch di chuyển vào phạm vi theo dõi của hệ thống quan sát điểm $A\left(5; 12; 0 \right)$.
	\par 
	\resizebox{!}{5cm}{\begin{tikzpicture}[line join = round,line cap=round,>=stealth]
		 \path
			(0,0) coordinate  (O)
			 (1.55,0) coordinate (C)
			(3.3,0) coordinate (D) 
			({3.3*cos (130)},{3.3*sin (130)}) coordinate (A) 
			({3.3*cos  (30)},{3.3*sin (30)}) coordinate (B)
			 ($(A)!1.5!(B)$) coordinate (B')
			 ($(B)!1.5!(A)$) coordinate (A')
			($(B)!1.3!(A)$) coordinate (M); 
		\fill[ball color=cyan,fill opacity=0.35] (O)  circle (1.55cm);
		\draw (O) circle  (3.3 cm);
		\draw[<->] 
			(O) -- (C)  node[pos=0.5,below]{$6370$ km};
		\draw[<->] 
			(D) -- (C) node[pos=0.5,above]{$6630$ km}; 
		\draw (A')  -- (B');
		 \fill
			(A)  circle (1pt) node[above]{$A$}
			(B) circle  (1pt) node[above]{$B$}
			(M) circle (1pt) node[above]{$M$}; 
	\end{tikzpicture}
	 }
	\choiceTF 
		{\True Đường thẳng $AM$ có phương trình chính tắc là~$\dfrac{x-5}{1} = \dfrac{y-12}{3}  = \dfrac{z}{-2}$}
		{Trên hệ trục tọa độ đã cho thiên thạch di chuyển qua điểm $N\left(7; 18; -5  \right)$}
		{\True Vị trí cuối cùng mà thiên thạch di chuyển trong phạm vi theo dõi của  hệ thống quan sát là~$B\left(-\dfrac{6}{7}; -\dfrac{39}{7}; \dfrac{82}{7} \right)$}
		{\True Khoảng cách giữa vị trí đầu tiên và vị trí cuối cùng ma thiên thạch di chuyển trong phạm vi theo dõi của hệ thống qua sát  là~$21\,915$ km (kết quả làm tròn đến hàng đơn vị theo đơn vị kilômét)}
	\loigiai{\begin{itemchoice}
	\itemch \textbf{Đúng}.\\ 
	Ta có $\overrightarrow{AM} = \left(1; 3; -2 \right)$ là~vectơ chỉ  phương của đường thẳng $AM$.\\
	Lại có $AM$ đi qua $A\left(5; 12;  0 \right)$ nên phương trình chính tắc của đường thẳng $AM$ là
	\begin{align*}
		 \dfrac{x-5}{1} = \dfrac{y-12}{3} = \dfrac{z}{-2}.
	\end{align*}
	\itemch \textbf{Sai}.\\ 
	Giả sử thiên thạch di chuyển qua điểm $N\left(7; 18; -5 \right)$  suy ra $N \in AM$.\\
	 Suy ra $\dfrac{7-5}{1} = \dfrac{18-12}{3}=\dfrac{-5}{-2} \Leftrightarrow 2 = 2 = \dfrac{5}{2}$ (vô lí).\\
	Vậy thiên thạch không di  chuyển qua điểm $N\left(7; 18; 5\right)$.
	\itemch  \textbf{Đúng}.\\
	Ta có $B \in AM \colon \dfrac{x-5}{1} = \dfrac{y-12}{3} = \dfrac{z}{-2}  \Rightarrow M\left(5+t; 12+3t; -2t \right)$.\\
	Ngoài thực tế khoảng cách từ tâm trái đất đến vị trí cuối cùng mà thiên thạch di chuyển trong phạm vi theo dõi của hệ thống quan sát là~$6370 + 6630 = 13\,000$ km. ứng với $13$ đơn vị trên trục toạ độ, hay \begin{align*}
	OB = 13 &\Leftrightarrow  OB^2 = 169\\
	&\Leftrightarrow \left(5+t \right)^2 + \left(12+3t  \right)^2 + \left(-2t \right)^2 = 169 \\
	&\Leftrightarrow 14t^2 + 82t =  0\\
	&\Leftrightarrow  \heva{&t=0\\&t=-\dfrac{41}{7}.}
\end{align*}
Với $t = 0$ ta được $B \left(5; 12; 0 \right) \equiv A$ (không  thoả mãn).\\
Với $t = -\dfrac{41}{7}$ ta được $B\left(-\dfrac{6}{7}; -\dfrac{39}{7}; \dfrac{82}{7}  \right)$ (thoả mãn).
 \itemch \textbf{Đúng}.\\
Ta có $AB = \sqrt{\left(-\dfrac{6}{7}-5 \right)^2 + \left(-\dfrac{39}{7} - 12  \right)^2 + \left(\dfrac{82}{7} \right)^2} = \dfrac{41\sqrt{14}}{7}$.\\
Khoảng  cách thực tế là~$1000 \cdot AB = 1000 \cdot \dfrac{41\sqrt{41}}{7} \approx 21\,915$ km.
\end{itemchoice}
 }
\end{ex}

% nguồn: Phan1_Lop_12_Hinh_hoc_Chuong_5
\begin{ex}%[ma:75b78fde]
	%[2H5C2-8] 
	Trong không gian với hệ tọa độ $Oxyz$ (đơn vị trên mỗi trục tọa độ là~kilômét), một máy bay đang ở vị trí $A\left( \dfrac {7}{2};-2;\dfrac {2}{5} \right)$ và sẽ hạ cánh ở vị trí $B\left( \dfrac {7}{2};\dfrac {11}{2};0 \right)$ trên đường băng $EG$. Biết rằng có một lớp mây được mô phỏng bởi một mặt phẳng $(\alpha)$ đi qua ba điểm $M\left(5;0;0 \right)$, $N\left(0;5;0\right)$, $P\left(0;0;\dfrac {1}{2} \right)$, điểm $C$ là~vị trí mà máy bay xuyên qua đám mây để hạ cánh và theo quy định an toàn bay, người phi công phải nhìn thấy điểm đầu $E\left(\dfrac {7}{2};\dfrac {9}{2};0 \right)$  của đường băng ở độ cao tối thiểu là~$120$ m (được mô phỏng bởi hình vẽ bên dưới). (\textit{Nguồn: R.Larson and B.Edwards, Calculus 10e, Cengage, 2014}).
	\begin{center}
		\resizebox{!}{5cm}{\begin{tikzpicture}[line join=round,line cap=round,>=stealth]
		\path (0:0) coordinate  (O)
			(-120:4)  coordinate (M)
			(-120:4.4) coordinate  (X)
			(90:2)  coordinate (P)
			(180:5)  coordinate (N)
			(180:5.2) coordinate  (N_1)
			(160:4) coordinate  (A)
			 ($(O)!1.5!(P)$) coordinate (Z)
			($(O)+(0:5)$) coordinate  (Y)
			(170:1.5) coordinate  (C)
			 ($(A)!1.5!(C)$) coordinate (D_1)
			($(A)!1.85!(C)$) coordinate  (D)
			($(A)!3!(C)$)  coordinate (B)
			($(B)+(0:2.2)$) coordinate  (G)
			($(B)+(180:2)$) coordinate (E); 
		\draw[->]  (O)--(Z) node[right]{$z$};
		\draw[->] (O)--(Y) node[below]{$y$}; 
		\draw[->] (O)--(X)  node[below right]{$x$};
		 \draw[dashed] (O)--(N)
			(C)--(D); 
		\draw (N)--(N_1) 
			(M)--(N)--(P)--cycle 
			(A)--(C) 
			(D_1)--(B) 
			 (E)--(B)--(G);
		\foreach \t/\g in  {M/0,N/90,P/180,A/90,C/90,B/-90,E/-90,G/-90}{
			\fill (\t) circle  (1pt) node[shift={(\g:7pt)},font=\scriptsize]{$ \t $};
		} 
		\path (P) node[above right,font=\scriptsize]{$\dfrac{1}{2}$} 
			(N) node[below  left,font=\scriptsize]{$-5$}
			(M) node[left,font=\scriptsize]{$5$}; 
	\end{tikzpicture}
		 }
	\end{center}
	 \choiceTF
		{\True Đường  thẳng $AB$ có phương trình tham số là~$\heva{& x=\dfrac {7}{2} \\ & y=-2+\dfrac {15}{2}t \\ & z=\dfrac {2}{5}-\dfrac {2}{5}t},\,( t\in \mathbb{R})$}
		{Tọa độ của điểm $C\left( \dfrac {7}{2};0;\dfrac {28}{115} \right)  $}
		{\True Khi máy bay đạt được độ cao $120$ m so với đường băng thì máy bay đang ở vị trí điểm $D$ trên đoạn  thẳng $AB$ có tọa độ là~$D\left( \dfrac {7}{2};\dfrac {13}{4};\dfrac {3}{25} \right) $}
		{Nếu tầm nhìn xa  của người phi công sau khi ra khỏi đám mây là~$900$ m thì người phi công đó đạt được quy định an toàn bay}
	\loigiai{ 
		 \par
		\resizebox{!}{5cm}{\begin{tikzpicture}[line join=round,line cap=round,>=stealth]
		\path (0:0) coordinate (O) 
			 (-120:4) coordinate (M)
			(-120:4.4)  coordinate (X)
			(90:2)  coordinate (P)
			(180:5) coordinate  (N)
			(180:5.2) coordinate (N_1) 
			(160:4)  coordinate (A)
			($(O)!1.5!(P)$) coordinate  (Z)
			 ($(O)+(0:5)$) coordinate (Y)
			 (170:1.5) coordinate (C)
			($(A)!1.5!(C)$) coordinate (D_1) 
			($(A)!1.85!(C)$) coordinate (D) 
			 ($(A)!3!(C)$) coordinate (B)
			($(B)+(0:2.2)$) coordinate  (G)
			 ($(B)+(180:2)$) coordinate (E);
		\draw[->] (O)--(Z)  node[right]{$z$};
		 \draw[->] (O)--(Y) node[below]{$y$};
		\draw[->] (O)--(X) node[below right]{$x$}; 
		\draw[dashed] (O)--(N) 
			(C)--(D); 
		 \draw (N)--(N_1)
			(M)--(N)--(P)--cycle 
			(A)--(C) 
			(D_1)--(B) 
			(D)--(E)--(B)--(G); 
		 \foreach \t/\g in {M/0,N/90,P/180,A/90,C/90,B/-90,D/90,E/-90,G/-90}{
			 \fill (\t) circle (1pt) node[shift={(\g:7pt)},font=\scriptsize]{$ \t $};
		} 
		\path (P)  node[right,font=\scriptsize]{$\dfrac{1}{2}$}
			(N) node[below  left,font=\scriptsize]{$-5$}
			(M)  node[left,font=\scriptsize]{$5$};
	\end{tikzpicture}
		} 
		\begin{itemchoice}
			 \itemch
			Đường thẳng $AB$ đi qua điểm $A\left( \dfrac {7}{2};-2;\dfrac {2}{5} \right)$ và nhận $\overrightarrow{AB}\left( 0;\dfrac {15}{2};-\dfrac {2}{5} \right)$ làm  véc-tơ chỉ phương.\\
			Phương trình tham  số của đường thẳng $AB$ là~$\heva{& x=\dfrac {7}{2} \\ & y=-2+\dfrac {15}{2}t,( t\in \mathbb{R})\cdot \\ & z=\dfrac {2}{5}-\dfrac {2}{5}t }$.\\
			 \itemch
			Mặt phẳng $( \alpha)$ đi qua ba điểm $M\left( 5;0;0  \right)$, $N\left( 0;5;0 \right)$, $P\left( 0;0;\dfrac {1}{2} \right)$ nên phương trình mặt phẳng $( \alpha)$ có dạng $\dfrac {x}{5}+\dfrac {y}{-5}+\dfrac {z}{\dfrac {1}{2}}=1$ hay $ x-y+10z-5=0 $.\\
			Vì $C$ là~vị trí mà máy bay xuyên qua đám  mây để hạ cánh nên $C$ là~giao điểm của đường thẳng $AB$ và mặt phẳng $( \alpha) $\\
			Vì $C\in AB$ nên có tọa độ  là~$C\left(\dfrac {7}{2};2+\dfrac {15}{2}t;\dfrac {2}{5}-\dfrac {2}{5}t \right) $\\
			$C\in  (\alpha)$ nên ta có $\dfrac {7}{2}-\left(-2+\dfrac {15}{2}t \right)+10\left( \dfrac {2}{5}-\dfrac {2}{5}t \right)-5=0\Leftrightarrow t=\dfrac {9}{23}$.\\
			Vậy $C\left( \dfrac {7}{2};\dfrac  {43}{46};\dfrac {28}{115} \right)$.\\
			 \itemch
			Vì $D\in AB$ nên  có tọa độ là~$D\left( \dfrac {7}{2};-2+\dfrac {15}{2}{t}';\dfrac {2}{5}-\dfrac {2}{5}{t}' \right) $\\
			$D$ là~vị trí mà tại đó máy bay ở độ cao $120$ m so với  đường băng, tức là~khoảng cách từ điểm $D$ đến mặt phẳng $( Oxy)$ bằng $120$ m $=0{,}12$ km.\\
			Ta có $\mathrm{d}\left( D,\left( Oxy \right) \right)=\left| \dfrac {2}{5}-\dfrac {2}{5}{t}' \right|=\dfrac {3}{25}\Leftrightarrow \hoac{& \dfrac {2}{5}-\dfrac  {2}{5}t'=\dfrac {3}{25} \\ & \dfrac {2}{5}-\dfrac {2}{5}t'=-\dfrac {3}{25} }\Leftrightarrow \hoac{& t'=\dfrac {7}{10} \\ & t'=\dfrac {13}{10}.} $\\
			Với $t'=\dfrac {7}{10}$, ta có $D\left( \dfrac {7}{2};\dfrac {13}{4};\dfrac {3}{25}  \right) $.\\
			Với $t'=\dfrac {13}{10},$ ta có $D\left( \dfrac {7}{2};\dfrac {31}{4};-\dfrac  {3}{25} \right) $.\\
			Vì $D$ là~vị trí độ cao của máy bay  nên ta chọn $D\left( \dfrac {7}{2};\dfrac {13}{4};\dfrac {3}{25} \right)$.\\
			\itemch 
			Ta có $DE=\sqrt {\left(\dfrac {7}{2}-\dfrac {7}{2} \right)^2+\left(  \dfrac {9}{2}-\dfrac {13}{4} \right)^2+\left( 0-\dfrac {3}{25} \right)^2}\approx 1{,}256$ (km).\\
			Vì tầm nhìn xa của phi công sau khi ra  khỏi đám mây là~$900$ m $=0{,}9$ km $<1{,}256$ km nên người phi công đó không đạt được quy định an toàn bay.
		\end{itemchoice}
	 }
\end{ex}

% nguồn: Phan1_Lop_12_Hinh_hoc_Chuong_5
\begin{ex}%[ma:904404f7]
	%[2H5H2-4] 
	Trong không  gian với hệ tọa~độ $Oxyz$, cho hai đường thẳng $\Delta\colon \dfrac{x-x_0}{a}=\dfrac{y-y_0}{b}=\dfrac{z-z_0}{c}$, $\Delta'\colon \dfrac{x-x'_0}{a'}=\dfrac{y-y'_0}{b'}=\dfrac{z-z'_0}{c'}$. Xét các điểm $M\left(x_0; y_0; z_0\right)$, $M'\left(x'_0; y'_0; z'_0\right)$ và các vectơ $\vec{u}=(a; b; c)$, $\overrightarrow{u'}=\left(a'; b'; c'\right)$.
	 \choiceTF
		{\True Nếu  hai vectơ $\vec{u}=(a; b; c)$, $\overrightarrow{u'}=\left(a'; b'; c'\right)$ cùng phương thì hai đường thẳng song song hoặc trùng nhau}
		{Nếu hai vectơ $\vec{u}=(a; b; c)$, $\overrightarrow{u'}=\left(a'; b'; c'\right)$ không cùng phương thì hai đường thẳng luôn  cắt nhau}
		{Nếu $\left[\vec{u}, \overrightarrow{u'}\right] \cdot \overrightarrow{MM'}=0$  thì hai đường thẳng luôn song song}
		{\True Nếu $\left[\vec{u}, \overrightarrow{u'}\right] \cdot \overrightarrow{MM'} \neq 0$ thì hai đường  thẳng chéo nhau}
	 \loigiai{
		\begin{itemchoice}
			\itemch \textbf{Đúng.} Vì  nếu hai vectơ $\vec{u}=(a; b; c)$, $\overrightarrow{u'}=\left(a'; b'; c'\right)$ cùng phương thì hai đường thẳng song song hoặc trùng nhau.
			\itemch \textbf{Sai.} Vì nếu hai vectơ $\vec{u}=(a; b; c)$, $\overrightarrow{u'}=\left(a'; b'; c'\right)$  không cùng phương thì hai đường thẳng chéo nhau hoặc cắt nhau.
			\itemch \textbf{Sai.}  Vì nếu $\left[\vec{u}, \overrightarrow{u'}\right] \cdot \overrightarrow{MM'}=0$ thì hai đường thẳng cắt nhau.
			\itemch \textbf{Đúng.} Vì nếu $\left[\vec{u}, \overrightarrow{u'}\right] \cdot  \overrightarrow{MM'} \neq 0$ thì hai đường thẳng chéo nhau.
		\end{itemchoice}
	} 
\end{ex}

% nguồn: Phan1_Lop_12_Hinh_hoc_Chuong_5
\begin{ex}%[ma:63fc5de7]
	 %[2H5H1-5]
	Trong  không gian $Oxyz$, cho hình lập phương $ABCD.A'B'C'D'$ có $A(0;0;0)$, $B(2;0;0)$, $D(0;2;2)$, $A'(0;0;2)$. Gọi $M$, $N$ lần lượt là~trung điểm của $AB$ và $AA'$ (Hình 3).
	\begin{center}
		\begin{tikzpicture}[>=stealth,line join=round,line cap=round]
				\def\a{3}; 
				\def\cao{3}; 
				 \def\g{30};
				 \path
					(0,0)  coordinate (B)
					(-1.5,-1) coordinate  (x)
					(1.3,5) coordinate  (z)
					(6,0.7) coordinate  (y)
					(\a,0) coordinate (C) 
					(\g:{\a/2})  coordinate (A)
					 ($(A)+(C)-(B)$) coordinate (D)
					($(B)+(90:\cao)$) coordinate  (B')
					($(C)+(90:\cao)$)  coordinate (C')
					($(D)+(90:\cao)$)  coordinate (D')
					($(A)+(90:\cao)$) coordinate  (A')
					(1.3,2.2) coordinate (N) 
					(0.65,0.35) coordinate (M); 
				\draw 
					(A')--(B')--(C')--(D')--cycle  (B)--(C)--(D)
					(C)--(C') 
					(B)--(B') 
					 (D)--(D');
				\draw[->] (B)--(x)  node[left]{$x$};
				\draw[->]  (A')--(z) node[left]{$z$};
				\draw[->] (D)--(y) node[above]{$y$}; 
				\draw[dashed] (B)--(A)--(D) 
					(A)--(A'); 
				 \draw[dashed] (M)--(N)--(D)--(M);
				 \foreach \x/\g in {A'/145,B'/180,C'/90,D'/90,A/140,B/-45,C/-45,D/-45,N/0,M/130}\fill (\x) circle (1pt) node[shift={(\g:10pt)}] {$\x$};
				\foreach \x/\y/\z in {}{\draw pic[draw,angle radius=2mm]{right angle=\x--\y--\z};} 
				\path (current bounding box.south) node[below=2mm] 
					{Hình  3};
			\end{tikzpicture}
	\end{center}
	\choiceTF 
		{\True  Tọa độ của điểm $M(1;0;0)$}
		{Tọa  độ của điểm $N(0;1;0)$}
		 {\True Phương trình mặt phẳng $(DMN)$ là~$\dfrac{x}{1}+\dfrac{y}{3}+\dfrac{z}{1}=1$}
		{\True Khoảng cách từ điểm $C'$ đến mặt  phẳng $(DMN)$ bằng $\dfrac{8}{3}$}
	 \loigiai{
		\begin{itemchoice}
			 \itemch Đúng. Do $A(0;0;0)$, $B(2;0;0)$ và $M$ là~trung điểm của $AB$ nên $M(1;0;0)$.
			\itemch Sai. Do $A(0;0;0)$, $A'(0;0;2)$ và $N$  là~trung điểm của $AA'$ nên $N(0;0;1)$.
			\itemch  Đúng. Phương trình mặt phẳng $(DMN)$ là~$\dfrac{x}{1}+\dfrac{y}{3}+\dfrac{z}{1}=1$ (phương tình mặt chắn).
			\itemch  Đúng. Ta có $\dfrac{x}{1}+\dfrac{y}{3}+\dfrac{z}{1}=1 \Leftrightarrow 2x+y+2z-2=0$. \\
			Mà điểm $C'(2;2;2)$ từ đó  ta có\\
			 $\mathrm{d}(C';(DMN))=\dfrac{|2\cdot2 +2+2\cdot 2-2|}{\sqrt{2^2+1^2+2^2}}=\dfrac{8}{3}$.
		\end{itemchoice}
	 }
\end{ex}

% nguồn: Phan1_Lop_12_Hinh_hoc_Chuong_5
\begin{ex}%[ma:3c4307e6]
	 %[2H5H2-7]
	Trong không gian $Oxyz$, cho mặt phẳng $(P)\colon x- y =0$. 
	 \choiceTF
		{\True  Mặt phẳng $(P)$ có một vectơ pháp tuyến~là $\vec{n}=(1; -1; 0)$}
		{\True Mặt phẳng $(P)$ vuông góc với mặt phẳng $(Oxy)$} 
		{Góc giữa mặt phẳng $(P)$  và trục $Oy$~là $135^\circ$}
		{\True Mặt  phẳng $(P)$ chứa trục $Oz$}
	 \loigiai{
		\begin{itemchoice}
			\itemch \textbf{Đúng.} Vì mặt phẳng $(P)\colon x- y  =0$ có một vectơ pháp tuyến~là $\vec{n}=(1; -1; 0)$.
			\itemch \textbf{Đúng.} Vì mặt phẳng $(Oxy)$ có một vectơ pháp  tuyến~là $\vec{k}=(0; 0; 1)$.\\
			Ta có $\vec{n}\cdot  \vec{k}=0\Rightarrow (P)\perp(Oxy)$.
			\itemch \textbf{Sai.} Vì trục $Oy$ có vectơ  đơn~vị $\vec{j}=(0;1;0)$.\\
			Ta có $\sin\left((P),Oy\right) =\dfrac{\left| \vec{n}\cdot\vec{j}\right| }{\left|\vec{n}\right|\cdot\left|\vec{j}\right|}=\dfrac{\left|-1\right|}{\sqrt{2}\cdot \sqrt{1}}=\dfrac{\sqrt{2}}{2}\Rightarrow \left((P),Oy\right)=45^\circ$. 
			\itemch \textbf{Đúng.}  Vì trục $Oz$ có vectơ đơn~vị $\vec{k}=(0;0;1)$.\\
			 Ta có $\heva{& \vec{n}\cdot \vec{k}=0 \\ & O\in (P)}\Rightarrow Oz \subset (P)$.
		\end{itemchoice}
	 }
\end{ex}

% nguồn: Phan1_Lop_12_Hinh_hoc_Chuong_5
\begin{ex}%[ma:8c0351fc]
	 %[2H5H2-7]%[2H5N1-2]%[2H5N2-2]
	Trong  không gian $Oxyz$, cho đường thẳng $d$ đi qua hai điểm $A(1; 2; 1)$ và $B(3; 0; 1)$, mặt phẳng $(\alpha)$ đi qua ba điểm $M(0; 1; 0)$, $N(2; 1; 3)$, $P(4; 1; 1)$.
	\choiceTF 
		{Véc-tơ $\overrightarrow{AB}$  không là~véc-tơ chỉ phương của đường thẳng $d$}
		{\True $\overrightarrow{MN} = (2; 0; 3)$,  $\overrightarrow{MP} = (4; 0; 1)$}
		{\True Mặt phẳng $(\alpha)$ có một véc-tơ pháp tuyến có tọa độ là~$(0; -1; 0)$} 
		{\True Góc giữa đường thẳng  $d$ và mặt phẳng $(\alpha)$ bằng $45^\circ$}
	\loigiai{\begin{itemchoice}
	\itemch Ta có  $\overrightarrow{AB} = (2; -2; 0)$ là~một véc-tơ chỉ phương của đường thẳng $d$.
	\itemch $\overrightarrow{MN} = (2;  0; 3)$, $\overrightarrow{MP} = (4; 0; 1)$.
	\itemch Ta có $\left[\overrightarrow{MN}, \overrightarrow{MP}\right] = (0; 10; 0)$. Mặt phẳng $(\alpha)$ có một véc-tơ  pháp tuyến là~$\overrightarrow{n} = (0; -1; 0)$.
	\itemch  Ta có
	 \[
	\sin(d, (\alpha)) = \dfrac{|\overrightarrow{AB} \cdot \overrightarrow{n}|}{|\overrightarrow{AB}| \cdot  |\overrightarrow{n}|} = \dfrac{2}{\sqrt{8} \cdot 1} = \dfrac{\sqrt{2}}{2}
	 \]
	Suy ra góc  giữa đường thẳng $d$ và mặt phẳng $(\alpha)$ bằng $45^\circ$.
\end{itemchoice}
} 
\end{ex}

% nguồn: Phan1_Lop_12_Hinh_hoc_Chuong_5
\begin{ex}%[ma:b690c95e]
	%[2H5V2-7]%[2H2V2-6] 
	Trong  không gian $Oxyz$, một cabin cáp treo xuất phát từ điểm $A(10;3;0)$ và chuyển động đều theo đường cáp có véc-tơ chỉ phương $\vec{u}=(2;-2;1)$ (hướng chuyển động cùng chiều với hướng véc-tơ $\vec{u}$ với tốc độ là~$4{,}5$ (m/s); đơn vị trên mỗi trục là~mét).
	\choiceTF 
		{\True Phương  trình tham số của đường cáp là~$\heva{&x=10+2k\\&y=3-2k\\&z=k,}(k\in\mathbb{R})$}
		{\True Giả sử sau thời gian $t$(s) kể  từ khi xuất phát ($t\geq 0$), cabin đến điểm $M$. Khi đó tọa độ điểm $M$ là~$\left(3t+10;-3t+3;\dfrac{3t}{2}\right)$}
		{Cabin dừng ở  điểm $B$ có hoành độ $x_B=550$, khi đó quãng đường $AB$ dài $800$ m}
		{Đường cáp $AB$ tạo với mặt phẳng $(Oxy)$ một  góc $30^\circ$}
	\loigiai{ 
		\begin{itemchoice}
			\itemch Phương trình tham số của đường thẳng $d$ qua $A(10;3;0)$ và có véc-tơ chỉ phương $\vec{u}=(2;-2; 1)$ là~$\heva{&x=10+2k\\&y=3-2k\\&z=k}(k\in\mathbb{R}).$ 
			 \itemch Ta có độ dài $AM=vt$.\\
			Vì $M$ thuộc đường thẳng $d$ nên  $M$ có tọa độ $(10+2k; 3-2k;k), (k\in\mathbb{R})$.\\
			 Vì $\vec{AM}=(2k;-2k;k)$ mà $\vec{AM}$ cùng hướng với $\vec{u}$ nên có $k \geq 0$, hay $AM=3k$.\\
			Vậy $3k=4{,}5t$ suy ra $k=1{,}5t$ do đó tọa độ điểm $M$  là~$\left(3t+10;-3t+3; \dfrac{3t}{2}\right)$.
			\itemch Ta thấy khi $x_B=550$ tức tà $3t+10=550$ thì $t=180$ (s).\\ 
			Vậy $AB=vt=4{,}5\cdot  180=810$ (m).
			\itemch Ta có $\vec{u}_{AB}=(2;-2;1)$. Ngoài ra, mặt phẳng ($Oxy$) có phương trình $z=0$ nên có một véc-tơ pháp tuyến $\vec{n}=(0;0;1)$.\\ 
			Gọi $\alpha$ là~góc giữa $AB$  và ($Oxy$), ta có $\sin \alpha=\dfrac{|\vec{u}\cdot \vec{n}|}{|\vec{u}|\cdot|\vec{n}|}=\dfrac{1}{3}$.\\
			Suy ra $\alpha\approx 19{,}47^\circ$. 
		\end{itemchoice}
	 }
\end{ex}

% nguồn: Phan1_Lop_12_Hinh_hoc_Chuong_5
\begin{ex}%[ma:2d9ea775]
	%[2H5H1-6] 
	Trong không gian $Oxyz$, cho hai mặt phẳng $(P)\colon 2x-y+z+1=0$ và  $(Q)\colon 4x-2y+2z+3=0$.
	\choiceTF 
		{$\vec{n}=(2;-1;-1)$~là một véc-tơ pháp  tuyến của $(P)$}
		 {\True $(P)\parallel(Q)$}
		 {Điểm $I\left(0;\dfrac{3}{2};0\right)$ không thuộc $(Q)$}
		{Khoảng cách  giữa hai mặt phẳng $(P)$ và $(Q)$ bằng $\dfrac{1}{\sqrt{6}}$}
	 \loigiai{
	\begin{itemchoice}
		\itemch \textbf{Sai}.  Ta có $\vec{n}=(2;-1;-1)$ không phải~là một véc-tơ pháp tuyến của~$(P)$.
		\itemch \textbf{Đúng}. Ta có $\dfrac{2}{4}=\dfrac{-1}{-2}=\dfrac{1}{2}\neq\dfrac{1}{3}$ nên  $(P)\parallel(Q)$.
		\itemch \textbf{Sai}. Ta có $4\cdot0-2\cdot\dfrac{3}{2}+2\cdot0+3=0$ nên điểm $I\left(0;\dfrac{3}{2};0\right)$ thuộc $(Q)$. 
		\itemch \textbf{Sai}. Khoảng cách giữa hai mặt phẳng  $(P)$ và $(Q)$~là
		 \[\mathrm{d}((P),(Q))=\mathrm{d}(I,(P))=\dfrac{\left|2\cdot0-\tfrac{3}{2}+0+1\right|}{\sqrt{4+1+1}}=\dfrac{1}{2\sqrt{6}}.\]
	\end{itemchoice}
	 }
\end{ex}

% nguồn: Phan1_Lop_12_Hinh_hoc_Chuong_5
\begin{ex}%[ma:5ea1c130]
	 %[2H5V1-6]
	Trong không gian $Oxyz$, $(P)$ là~mặt phẳng qua  $M(1;2;3)$ và song song với $(Q)\colon x-y+2z=0$.
	 \choiceTF
		{\True Một vectơ pháp tuyến  của mặt phẳng $(P)$ là~$\overrightarrow{n}=(1;-1;2)$}
		 {Phương trình mặt phẳng $(P)$ là~$x-y+2z-4=0$}
		{Khoảng  cách từ gốc tọa độ $O$ đến mặt phẳng $(P)$ bằng $\dfrac{6}{7}$}
		{\True Côsin góc giữa mặt phẳng  $(P)$ và mặt phẳng $(Oxz)$ bằng $\dfrac{\sqrt{6}}{6}$}
	\loigiai{ 
		\begin{itemchoice}
			\itemch $(P)$ song song với $(Q):x-y+2z=0$ mà $(Q)$ có véc-tơ pháp tuyến là~$\overrightarrow{n}=(1;-1;2)$ nên $(P)$ cũng nhận  $\overrightarrow{n}$ làm véc-tơ pháp tuyến.
			\itemch $(P)$ qua $M(1;2;3)$, có véc-tơ pháp tuyến là~$\overrightarrow{n}=(1;-1;2)$  nên $(P):x-y+2z-5=0$.
			\itemch Ta có $\mathrm{d}\left( O, (P) \right) = \dfrac{|-5|}{\sqrt{1^2+(-1)^2+2^2}}  = \dfrac{5\sqrt{6}}{6}$.
			 \itemch $(Oxz)$ có véc-tơ pháp tuyến là~$\overrightarrow{j}=(0;1;0)$.
			$$\cos\left( (P),(Oxz) \right) =  \dfrac{\left| \overrightarrow{n}\cdot \overrightarrow{j} \right|}{\left| \overrightarrow{n} \right| \cdot \left| \overrightarrow{j} \right|}
			= \dfrac{\left| 1\cdot 0 - 1\cdot 1  + 2\cdot 0 \right|}{\sqrt{1^2+(-1)^2+2^2} \cdot \sqrt{0^2+1^2+0^2}}
			 = \dfrac{1}{\sqrt{6}}.$$
		\end{itemchoice}
	} 
\end{ex}

% nguồn: Phan1_Lop_12_Hinh_hoc_Chuong_5
\begin{ex}%[ma:cf9ee185]
	 %[2H5H3-4]
	Trong không gian $Oxyz$ (đơn~vị trên mỗi trục tính theo kilômét), một trạm thu  phát sóng điện thoại di động được đặt ở vị trí $I(1;3;7)$. Trạm thu phát sóng đó được thiết kế với bán kính phủ sóng~là $3 \mathrm{~km}$.
	 \choiceTF
		{Phương trình mặt cầu $(S)$ để mô tả  ranh giới bên ngoài của vùng phủ sóng trong không gian~là $(x+1)^2+(y+3)^2+(z+7)^2=9$}
		 {Điểm $A(2;2;7)$ nằm ngoài mặt cầu $(S)$}
		{\True Nếu người  dùng điện thoại ở vị trí có toạ độ $(2;2;7)$ thì có thể sử dụng dịch vụ của trạm thu phát sóng đó}
		{\True Nếu người dùng điện thoại ở vị trí có toạ độ $(5;6;7)$ thì không thể sử dụng dịch vụ của  trạm thu phát sóng đó}
	 \loigiai{
		\begin{itemchoice}
			\itemch \textbf{Sai}.\\ Phương trình mặt cầu $(S)$  để mô tả ranh giới bên ngoài của vùng phủ sóng trong không gian~là $(x-1)^2+(y-3)^2+(z-7)^2=9$.
			\itemch \textbf{Sai}.\\ Ta có  $IA=\sqrt{2}<3$.\\
			Suy ra  điểm $A(2;2;7)$ nằm trong mặt cầu $(S)$.
			\itemch \textbf{Đúng}.\\ Ta có  $IA=\sqrt{2}<3$.\\
			Suy ra điểm $A(2;2;7)$ nằm  trong mặt cầu $(S)$.\\
			Vậy người dùng  điện thoại ở vị trí có toạ độ $(2;2;7)$ thì có thể sử dụng dịch vụ của trạm thu phát sóng đó.
			\itemch \textbf{Đúng}.\\ Gọi $M(5;6;7)$~là vị trí của người dùng  điện thoại. Ta có\\
			 $IM=5>3$.\\
			Vậy người dùng điện thoại ở vị trí có toạ độ $(5;6;7)$ không thể sử dụng dịch  vụ của trạm thu phát sóng này.
		\end{itemchoice}
	} 
\end{ex}

% nguồn: Phan1_Lop_12_Hinh_hoc_Chuong_5
\begin{ex}%[ma:adefeecc]
	%[2H5H2-4] 
	Trong không gian với hệ  tọa~độ $Oxyz$, cho hai đường thẳng
	\begin{align*}
		\Delta\colon \dfrac{x-1}{1}=\dfrac{y+1}{2}=\dfrac{z+3}{3},\;\Delta'\colon \dfrac{x}{3}=\dfrac{y-1}{2}=\dfrac{z-2}{1}\cdot 
	\end{align*}
	 \choiceTF
		{\True Hai vectơ $\vec{u}=(1; 2; 3),$ $ \vec{u}^\prime=(3; 2; 1)$ lần lượt~là hai vectơ chỉ phương  của $\Delta,$ $ \Delta'$}
		{Điểm $M(1;-1;-3)$ không thuộc đường thẳng $\Delta$, điểm $M'(0; 1;  2)$ không thuộc đường thẳng $\Delta'$}
		{\True $\left[\vec{u},  \vec{u}'\right]=(-4; 8;-4)$}
		{Hai  đường thẳng $\Delta,$ $ \Delta'$ chéo nhau}
	 \loigiai{
		\begin{itemchoice}
			\itemch  \textbf{Đúng.} Vì $\Delta\colon \dfrac{x-1}{1}=\dfrac{y+1}{2}=\dfrac{z+3}{3} $ nên nhận $\vec{u}=(1; 2; 3)$~làm vectơ chỉ phương, $\Delta'\colon \dfrac{x}{3}=\dfrac{y-1}{2}=\dfrac{z-2}{1}$ nên nhận $\vec{u}^\prime=(3; 2; 1)$~làm vectơ chỉ phương.
			\itemch \textbf{Sai.} Vì $\Delta\colon \dfrac{x-1}{1}=\dfrac{y+1}{2}=\dfrac{z+3}{3} $ nên  $\Delta$ đi qua $M(1;-1;-3)$. \\
			$\Delta'\colon  \dfrac{x}{3}=\dfrac{y-1}{2}=\dfrac{z-2}{1}$ nên $\Delta'$ đi qua $M'(0; 1; 2)$.
			\itemch \textbf{Đúng.} Vì 
			 $$\vec{u}=(1; 2; 3), \vec{u}^\prime=(3; 2; 1)\Rightarrow \left[\vec{u}, \vec{u}'\right]=(-4; 8;-4).$$
			\itemch \textbf{Sai.} Vì ta có 
			\begin{itemize}
				\item $\Delta$ đi  qua $M(1;-1;-3)$, có vectơ chỉ phương $\vec{u}=(1; 2; 3).$
				\item $\Delta'$ đi qua $M'(0; 1; 2)$,  có vectơ chỉ phương $\vec{u}^\prime=(3; 2; 1)$.
			\end{itemize}
			 Khi đó
			$$\overrightarrow{MM'}=(-1;2;5)\Rightarrow \left[\vec{u}, \vec{u}'\right]\cdot \overrightarrow{MM'}=-4\cdot (-1)+8\cdot 2+(-4)\cdot 5=0.$$  nên $\Delta,$ $ \Delta'$ không chéo nhau.
		\end{itemchoice}
	 }
\end{ex}

% nguồn: Phan1_Lop_12_Hinh_hoc_Chuong_5
\begin{ex}%[ma:8f189147]
	%[2H5H1-6] 
	Trong không gian với hệ tọa~độ $Oxyz$, cho hai mặt phẳng $\left(P_1\right)\colon 3x+ 4y+ 12z + 2025 = 0$, $\left(P_2\right)\colon 4x- 3y + 12z -  2026 = 0$.
	\choiceTF 
		{\True $\vec{n}_1=(3; 4; 12)$~là  một vectơ pháp tuyến của $\left(P_1\right)$}
		 {\True $\vec{n}_2=(4; -3; 12)$~là một vectơ pháp tuyến của $\left(P_2\right)$}
		{\True  $\vec{n}_1 \cdot \vec{n}_2=144$}
		 {\True Góc giữa hai mặt phẳng đã cho (làm tròn đến hàng đơn~vị của~độ)~là $32^\circ$}
	 \loigiai{
		\begin{itemchoice}
			\itemch \textbf{Đúng.} Vì mặt phẳng $\left(P_1\right)\colon 3x+ 4y+ 12z + 2025 = 0$ có một vectơ pháp tuyến $\vec{n}_1=(3;  4; 12)$.
			\itemch  \textbf{Đúng.} Vì mặt phẳng $\left(P_2\right)\colon 4x- 3y + 12z - 2026 = 0$ có một vectơ pháp tuyến $\vec{n}_2=(4; -3; 12)$.
			\itemch \textbf{Đúng.} Vì  $\vec{n}_1\cdot\vec{n}_2=3\cdot4+4\cdot(-3)+12\cdot12=144$.
			\itemch \textbf{Đúng.} Vì 
			 \allowdisplaybreaks
			\begin{eqnarray*}
				\cos\left((P_1),(P_2)\right)&=&\dfrac{\left| \vec{n}_1\cdot\vec{n}_2\right| }{\left|\vec{n}_1\right|\cdot\left|\vec{n}_2\right|}\\ 
				&=&\dfrac{\left| 144\right|}{\sqrt{3^2+4^2+12^2}\cdot\sqrt{4^2+(-3)^2+12^2}}=\dfrac{144}{169}\\  &\Rightarrow& \left((P_1),(P_2)\right)\approx 32^\circ
			\end{eqnarray*}
		\end{itemchoice}
	 }
\end{ex}

% nguồn: Phan1_Lop_12_Hinh_hoc_Chuong_5
\begin{ex}%[ma:723f7e56]
	%[2H5H3-4] 
	Trong không gian $Oxyz$, (đơn vị trên mỗi trục tính theo mét), một ngọn hải đăng được đặt ở vị trí $  I\left(17;20;45\right)$. Biết rằng ngọn hải đăng đó được thiết kế với bán kính phủ sáng là~$ 4$ km.
	\choiceTF 
		{\True Phương trình mặt cầu để mô tả ranh giới bên  ngoài của vùng phủ sáng trên biển của hải đăng là~$(x-17)^2+(y-20)^2+(z-45)^2=16\ 000\ 000$}
		{Nếu người  đi biển ở vị trí $ M\left(18;21;50\right)$ thì không thể nhìn thấy được ánh sáng từ ngọn hải đăng}
		{\True Nếu người đi biển ở vị trí $ N(4019;21;44)$ thì có thể nhìn thấy được ánh  sáng từ ngọn hải đăng}
		{\True Nếu hai người đi biển ở vị trí có thể nhìn thấy được ánh  sáng từ ngọn hải đăng thì khoảng cách giữa hai người đó không quá $ 8$ km}
	\loigiai{ 
		\begin{itemchoice}
			\itemch Phương trình mặt cầu tâm $  I(17;20;45)$ bán kính $ R=4000$ m là~$$(x-17)^2+(y-20)^2+(z-45)^2=16\ 000\ 000.$$
			\itemch  $ IM=\sqrt{(18-17)^2+(21-20)^2+(50-45)^2}=\sqrt{27}<16000000$.\\
			Suy ra người ở vị  trí điểm $M$ vẫn nhìn thấy ánh sáng từ ngọn hải đăng.
			\itemch $ IN=\sqrt{(4019-17)^2+(21-20)^2+(44-45)^2}=\sqrt{16016006}<16000000$. Suy ra người ở vị trí điểm $ N$ vẫn nhìn thấy  ánh sáng từ ngọn hải đăng.
			\itemch Vì đuờng kính của mặt cầu trên bằng $ 8000$ m hay $8$ km nên hai người đi biển ở vị trí có thể nhìn thấy được ánh sáng từ ngọn hải đăng thì khoảng cách giữa hai người  đó không quá $ 8$ km.
		\end{itemchoice}
	} 
\end{ex}

% nguồn: Phan1_Lop_12_Hinh_hoc_Chuong_5
\begin{ex}%[ma:e829b376]
	 %[2H5H1-3]
	Trong không gian $Oxyz$, cho đường thẳng $d\colon  \dfrac{x+2}{1}=\dfrac{y-1}{2}=\dfrac{z}{-2}$ và điểm $A\left(-3;-1;3\right)$ . Mặt phẳng $\left(P\right)$ đi qua $A$ và chứa đường thẳng $d$.
	 \choiceTF
		{\True Một vectơ chỉ phương  của $d$ là~$\overrightarrow{a}=\left(1;2;-2\right)$}
		{Đường thẳng $d$ qua  điểm $M\left(3;11;-9\right)$}
		{\True Phương trình tham số của $d$ có dạng $\heva{&x=-2-t\\&y=1-2t\\&z=2t}$} 
		{Phương trình mặt phẳng  $\left(P\right)$ là~$x+2y-2z+11=0$}
	 \loigiai{
		\begin{itemchoice}
			\itemch Một vectơ chỉ phương của $d$ là~$\overrightarrow{a}=(1;2;-2)$. 
			\itemch Đường  thẳng $d$ không qua điểm $M(3;11;-9)$ vì $\dfrac{3+2}{1}=\dfrac{11-1}{2}\ne\dfrac{-9}{-2}$.
			\itemch $d$ có một vectơ chỉ phương là~$\overrightarrow{b}=(-1;-2;2)$ và qua  $N(-2;1;0)$ nên $d$ có phương trình tham số là~$\heva{&x=-2-t\\&y=1-2t\\&z=2t.}$
			\itemch Ta có  $\overrightarrow{AN}=(1;2;-3)$.\\
			 Mặt phẳng $(P)$ có 1 vectơ pháp tuyến là~$\overrightarrow{n}=\Big[\overrightarrow{AN},\overrightarrow{a}\Big]=(2;-1;0)$.\\
			Phương  trình $(P)$ có dạng $2(x+2)-1(y-1)=0\Leftrightarrow 2x-y+5=0$.
		\end{itemchoice}
	 }
\end{ex}

% nguồn: Phan1_Lop_12_Hinh_hoc_Chuong_5
\begin{ex}%[ma:42604021]
	 %[2H5C1-7]
	Trong một cuộc thi thể thao về môn bắn súng, các vận động viên phải thực hiện bắn hạ mục tiêu đang di động trên mặt đất của khối cầu đặc có bán kính bằng $1$ m. Chọn hệ trục toạ độ Oxyz trong không gian có gốc $O$ đặt tại vị trí xạ thủ A ngắm bắn, xem mặt phẳng $(Oxy)$ là~mặt đất, đơn vị độ dài trên mỗi trục  tọa độ là~$1$ m. Biết khối cầu có tâm $I(7;24;3)$ và xem đường đi của viên đạn là~một đường thẳng.
	\choiceTF 
		{\True Vị trí xa nhất để xạ thủ  A nhìn thấy và ngắm bắn mục tiêu là~$25{,}2$ m (làm tròn đến hàng phần mười)}
		{\True Biết vận tốc viên đạn là~$\dfrac{54}{5}\sqrt{65} \ (\text{km/h})$ thì khoảng thời gian ngắn nhất để  xạ thủ A bắn trúng mục tiêu chưa tới $1$ s}
		{\True Để các xạ thủ có thể dễ dàng bắn trúng mục tiêu hơn, ban tổ chức đã quyết định cho mục tiêu di chuyển trên đường tròn lớn nhất của mặt cầu và song song với mặt  đất. Khi đó khoảng cách ngắn nhất từ vị trí xạ thủ A ngắm bắn đến mục tiêu là~$ 3\sqrt{65} \ (\text{m})$}
		{\True Xạ thủ A đang ngắm ở vị trí gần mục tiêu nhất. Tại thời điểm tuyển thủ A nổ súng thì mục tiêu đang ở vị trí $M(6;24;3)$ di  chuyển với vận tốc $v = \arctan\left(\dfrac{24}{7}\right) \ (\text{m/s})$ và đi ngược chiều kim đồng hồ. Khi đó xạ thủ A bắn trúng mục tiêu}
	\loigiai{ 
		\begin{itemchoice}
			\itemch  \textbf{Đúng}
			Gọi $M$ là~điểm mà tại đó mục  tiêu ở vị trí xa nhất.\\
			 Ta có
			$$ 
			AM = OM = \sqrt{OI^2 - IM^2} = \sqrt{7^2 + 24^2 +  3^2 - 1^2} = \sqrt{633} \approx 25{,}2\ \text{ (m)}.
			 $$
			\itemch \textbf{Đúng} Gọi $N$ là~điểm mà tại đó  mục tiêu ở vị trí gần nhất.\\
			Ta có: 
			$$ 
			AN = AI - IN = OI  - R = \sqrt{7^2 + 24^2 + 3^2} - 1 = \sqrt{634} - 1\ \text{ (m)}.
			 $$
			Vậy thời gian  ngắn nhất để xạ thủ $A$ bắn trúng mục tiêu là
			$$ 
			t = \dfrac{AN}{v} =  \dfrac{\sqrt{634} - 1}{\dfrac{54}{5} \sqrt{65}} \approx 0{,}9\,997\ \text{ (s)}.
			 $$
			\itemch  \textbf{Đúng}
			Mặt phẳng $(\alpha)$ đi qua điểm $I(7; 24; 3)$ và song song với mặt đất $(Oxy)$ nên có vectơ pháp tuyến $\vec{n}_{\alpha} = (0; 0;  1)$.\\
			Phương trình mặt phẳng $(\alpha)$ là~$0(x - 7) + 0(y - 24) + 1(z - 3) = 0 \Leftrightarrow z = 3$.\\ 
			Gọi $P$ là~vị trí mà tại đó mục  tiêu ở vị trí gần nhất. Do $P \in (\alpha)$ nên ta có $P(x_p; y_p; 3)$.\\
			Vì  $IP = R = 1$ nên
			\begin{eqnarray*}
				& & IP^2  = R^2 = 1 \\
				\Rightarrow & &(x_p - 7)^2  + (y_p - 24)^2 = 1 \\
				\Rightarrow & &x_p^2 - 14x_p +  49 + y_p^2 - 48y_p + 576 = 1 \\
				\Rightarrow & &x_p^2 +  y_p^2 = 14x_p + 48y_p - 624. \qquad(1)
			\end{eqnarray*}
			 Áp dụng bất đẳng thức Bunhiacopxki
			 $$
			x_p^2 + y_p^2 = 14x_p + 48y_p - 624 \leq \sqrt{14^2 + 48^2}  \cdot \sqrt{x_p^2 + y_p^2} - 624.
			$$ 
			 Gọi $t = x_p^2 + y_p^2$, ta có
			\begin{align*}
				t  &\leq 50\sqrt{t} - 624 \\
				\Rightarrow t + 624 &\leq  50\sqrt{t} \\
				 \Rightarrow (t + 624)^2 &\leq 2\,500t \\
				\Rightarrow t^2 +  1\,248t + 389\,376 &\leq 2\,500t \\
				\Rightarrow t^2 - 1\,252t  + 389\,376 &\leq 0
			\end{align*}
			Giải bất  phương trình ta được
			\begin{align*}
				 576 \leq t \leq 676 \\
				\Rightarrow 576 \leq  x_p^2 + y_p^2 \leq 676.
			\end{align*}
			Khi $x_p^2  + y_p^2 = 576$ thì $OP = \sqrt{x_p^2 + y_p^2 + 3^2} = \sqrt{576 + 9} = \sqrt{585}$.\\
			Tuy nhiên đề bài cho khoảng cách ngắn nhất là~$$OP = \sqrt{576 + 3^2} = \sqrt{576 + 9} = \sqrt{585} = 3\sqrt{65}  \ (\text{m}).$$
			Vậy khoảng cách ngắn nhất từ  xạ thủ $A$ đến mục tiêu là~$3\sqrt{65}$ (m).
			\itemch \textbf{Đúng} 
			Dấu  xảy ra (tức xạ thủ A bắn trúng mục tiêu) khi và chỉ khi
			$$ 
			 \dfrac{x_p}{14} = \dfrac{y_p}{48} \Leftrightarrow \dfrac{1}{14^2} x_p^2 - \dfrac{1}{48^2} y_p^2 = 0.
			 $$
			 Ta có hệ phương trình
			$$ 
			\begin{cases}
				\dfrac{1}{14^2} x_p^2 -  \dfrac{1}{48^2} y_p^2 = 0 \\
				x_p^2 +  y_p^2 = 576
			\end{cases}
			 \Rightarrow
			\begin{cases}
				x_p^2 = \dfrac{28\,224}{625}  \Rightarrow x_p = 6{,}72 \\
				y_p^2 = 530{,}8\,416 \Rightarrow y_p =  23{,}04.
			\end{cases}
			$$ 
			Khi đó xạ thủ A đang ngắm  ở điểm $P(6{,}72;\ 23{,}04;\ 3)$.\\
			Gọi  $H$ là~trung điểm của đoạn $MP$, khi đó
			 $$
			H\left( \dfrac{6{,}72 + 6}{2},\ \dfrac{23{,}04 +  24}{2},\ \dfrac{3 + 3}{2} \right) = H(6{,}36;\ 23{,}52;\ 3).
			$$ 
			Tính khoảng cách $IH = \sqrt{(6{,}36 - 7)^2 + (23{,}52 - 24)^2 + (3 -  3)^2} = 0{,}8$.\\
			Tính góc  $\widehat{HIM}$
			\begin{align*}
				\tan \widehat{HIM} &= \dfrac{HM}{IH} = \dfrac{\sqrt{(6 - 6{,}36)^2 + (24 - 23{,}52)^2 + (3 - 3)^2}}{0{,}8} =  0{,}75 \\
				\tan \widehat{MIP} &= \tan(2 \cdot \widehat{HIM}) = \dfrac{2 \cdot 0{,}75}{1 - 0{,}75^2} =  \dfrac{2{,}0 \cdot 0{,}75}{1 - 0{,}5\,625} = \dfrac{1{,}5}{0{,}4\,375} = \dfrac{24}{7}
			\end{align*}
			Suy ra 
			 $$
			\widehat{HIP}  = \arctan\left(\dfrac{24}{7}\right).
			 $$
			Thời gian mục tiêu di  chuyển từ $M$ đến $P$ là
			$$ 
			 \dfrac{\arctan\left(\dfrac{24}{7}\right)}{\arctan\left(\dfrac{24}{7}\right)} = 1\ \text{ (s)}
			 $$
			Thời gian viên đạn đi từ $O$ đến $P$ là 
			 $$
			\dfrac{3\sqrt{65}}{\dfrac{54}{5} \cdot \sqrt{65}} = \dfrac{3}{\dfrac{54}{5}} = \dfrac{15}{54} = \dfrac{5}{18} \cdot 3{,}6  = 1\ \text{ (s)}.
			 $$
			Vậy xạ thủ A  bắn trúng mục tiêu.
		\end{itemchoice}
	 }
\end{ex}

% nguồn: Phan1_Lop_12_Hinh_hoc_Chuong_5
\begin{ex}%[ma:c1a89582]
	 %[2H5N2-2]%[2H5N1-2]%[2H5H2-7]
	Trong  không gian $Oxyz$, cho đường thẳng $\Delta\colon\dfrac{x-2\,024}{2}=\dfrac{y}{1}=\dfrac{z+2\,025}{-2}$
	và mặt phẳng $(P)\colon 2x+2y-z+1=0$.  Xét các vectơ $\vec{u}=(2;1;-2)$, $\vec{n}=(2;2;-1)$.
	\choiceTF 
		 {\True $\vec{u}$~là một vectơ chỉ phương của đường thẳng $\Delta$}
		{\True $\vec{n}$~là một vectơ pháp tuyến của mặt phẳng $(P)$} 
		 {$\cos\big(\Delta,(P)\big)=\dfrac{8}{9}$}
		{\True Góc giữa đường thẳng $\Delta$ và mặt phẳng $(P)$ bằng $63^\circ$ (làm tròn đến  hàng đơn~vị của~độ)}
	\loigiai{ 
		\begin{itemchoice}
			 \itemch \textbf{Đúng}.\\
			Đường thẳng $\Delta\colon\dfrac{x-2\,024}{2}=\dfrac{y}{1}=\dfrac{z+2\,025}{-2}$ có một vectơ  chỉ phương~là $\vec{u}=(2;1;-2)$.
			\itemch  \textbf{Đúng}.\\
			Mặt phẳng $(P)\colon 2x+2y-z+1=0$ có một vectơ pháp tuyến~là $\vec{n}=(2;2;-1)$. 
			\itemch \textbf{Sai}.\\ 
			 $\sin\left(\Delta,(P)\right) =\dfrac{|\vec{u}\cdot\vec{n}|}{|\vec{u}|\cdot|\vec{n}|} =\dfrac{|2\cdot2+1\cdot2+(-2)\cdot(-1)|}{\sqrt{2^2+1^2+(-2)^2}\cdot\sqrt{2^2+2^2+(-1)^2}} =\dfrac{8}{9}$.\\
			$\cos\left(\Delta,(P)\right)=\sqrt{1-\left(\dfrac{8}{9}\right)^2}=\dfrac{\sqrt{17}}{9}$. 
			\itemch \textbf{Đúng}.\\ 
			Ta có $\sin\left(\Delta,(P)\right)=\dfrac{8}{9}  \Rightarrow \left(\Delta,(P)\right)=62^\circ44'2''$.
		\end{itemchoice}
	 }
\end{ex}

% nguồn: Phan1_Lop_12_Hinh_hoc_Chuong_5
\begin{ex}%[ma:191fb627]
	 %[2H5C3-4]
	Trong một cuộc thi thể thao về môn bắn súng. Các vận động viên phải thực hiện bắn hạ mục tiêu đang di động trên mặt của  khối cầu đặc có bán kính bằng $1$ m. Chọn hệ trục tọa độ $(Oxyz)$ trong không gian có gốc $O$ đặt tại vị trí xạ thủ $A$ ngắm bắn, xem mặt phẳng $(Oxy)$ là~mặt đất, đơn vị độ dài trên mỗi trục tọa độ là~$1$ m. Biết khối cầu có tâm $I(7;24;3)$ và xem đường đi của viên đạn là~một đường thẳng.
	 \choiceTF
		{\True Vị trí xa nhất để xạ thủ $A$ nhìn thấy và ngắm bắn mục tiêu là~$25,2$ m (làm tròn  đến hàng phần mười)}
		{\True Biết vận tốc viên đạn là~$\dfrac{54}{5}\sqrt{65}$ km/h thì  khoảng thời gian ngắn nhất để xạ thủ $A$ bắn trúng mục tiêu chưa tới $1$ giây}
		{\True Để các xạ thủ có thể dễ dàng bắn trúng mục tiêu hơn, ban tổ chức đã quyết định cho mục tiêu di chuyển trên đường tròn lớn nhất của mặt cầu  và song song với mặt đất. Khi đó khoảng cách ngắn nhất từ vị trí xạ thủ $A$ ngắm bắn đến mục tiêu là~$3\sqrt{65}$ m}
		{\True Xạ thủ $A$ đang ngắm ở vị trí gần mục tiêu nhất. Tại thời điểm tuyển thủ $A$ nổ súng thì mục tiêu đang ở vị trí $M(6;24;3)$ di chuyển với  vận tốc $v=\arctan\left(\dfrac{24}{7}\right)$ (m/s) và đi ngược chiều kim đồng hồ. Khi đó xạ thủ $A$ bắn trúng mục tiêu}
	\loigiai{ 
		\par 
		\resizebox{!}{6cm}{\begin{tikzpicture}[line join=round,line cap=round]
		 \path
			(3,10)  coordinate (I)
			 ($(I)-(3,0)$) coordinate (M)
			($(I)-(0.3,0.8)$)  coordinate (N)
			 ($(I)+(2,2.25)$) coordinate (E)
			($(I)-(2.12,2.12)$) coordinate  (F)
			($(I)-(-1.3,2.7)$) coordinate  (T)
			(3,1)  coordinate (H)
			 (-5,9) coordinate (K)
			(-6,3)  coordinate (A)
			(11,3) coordinate (B) 
			 (-11,0) coordinate (D)
			($(B)+(D)-(A)$) coordinate (C) 
			(-5,2) coordinate  (O);
		 \draw[dashed,thin]
			(M) arc (180:0:3  cm and 0.8 cm);
		\draw 
			(I) circle (3 cm) 
			 (M) arc (-180:0:3 cm and 0.8 cm);
		\draw  (A)--(B)--(C)--(D)--cycle;
		\draw[->] (O)--(-8,0.3) node[left]{$  x $};
		\draw[->] (O)--(-1,2) node[below]{$ y  $};
		 \draw
			 (I)--(O)
			(N)--(O) 
			 (T)--(O)
			(K)--(O); 
		\draw[dashed] 
			(I)--(K) 
			 (I)--(H)
			(I)--(T) 
			(I)--(E); 
		\node  at ($(I)+(1.1,0.05)$) {$(7,24,3)$};
		\node at ($(H)+(0.8,-0.35)$)  {$(7,24,0)$};
		 \node at ($(K)+(0.7,0.3)$) {$(0,0,3)$};
		\foreach \i/\g in  {O/-90,I/0,N/-90,F/180,T/0,K/90,H/-90}{\fill (\i) circle (1pt)
			 ($(\i)+(\g:3.3mm)$) node[scale=1]{$\i$};}
	\end{tikzpicture}
		 }
		\begin{itemchoice}
			\itemch Điểm xa nhất mà xạ thủ $A$ thấy được là~tiếp điểm $B$ của tiếp tuyến kẻ  từ $O$ đến mặt cầu.\\
			Ta có $OB^2=OI^2-IB^2=7^2+{{24}^2}+3^2-1^2\Rightarrow OB=\sqrt{633}\approx  25{,}2$ (m).
			\itemch Vì vận tốc không đổi nên khoảng thời gian ngắn nhất để xạ thủ $A$ bắn trúng mục tiêu là~khoảng thời gian cho quãng đường từ xạ thủ đến vị trí  gần xạ thủ nhất.\\
			Ta có $OC=OI-R=\sqrt{7^2+{{24}^2}+3^2}-1=\sqrt{634}-1$ (m).\\ 
			$v=\dfrac{54}{5}\sqrt{65}$ (km/h) $=3\sqrt{65}$  (m/s).\\
			Mặt khác, $t=\dfrac{s}{v}=\dfrac{OC}{v}=\dfrac{\sqrt{634}-1}{3\sqrt{65}}\approx 0{,}99969<1$ giây. 
			\itemch 
			Gọi $H(7;24;0)$ là~hình chiếu của $I$ lên mặt phẳng $(Oxy)$. Vì đường tròn lớn nhất của mặt cầu nằm trong mặt phẳng song song với mặt đất nên khoảng cách ngắn nhất là~$OD$ với $D$ là~một trong  hai giao điểm của mặt cầu, mặt phẳng $z=3$ và mặt phẳng $(OIH)$.\\
			 Ta có $ID=1$; $OI=\sqrt{634}-1$; $\cos \widehat{DIO}=\cos \widehat{IOH}=\cos (\overrightarrow{OH},\overrightarrow{OI})=\dfrac{\sqrt{7^2+{{24}^2}}}{\sqrt{7^2+{{24}^2}+3^2}}=\dfrac{25}{\sqrt{634}}$.\\
			Suy ra $OD=\sqrt{OI^2+ID^2-2OI\cdot  ID\cdot\cos \widehat{DOI}}=\sqrt{634+1-2\sqrt{634}.1.\dfrac{25}{\sqrt{634}}}=3\sqrt{65}$ (m).
			\itemch 
			 Ta có $\overrightarrow{IM}=(-1;0;0)$, $\overrightarrow{IO}=(-7;-24;-3)\Rightarrow \widehat{MIC}=(\overrightarrow{IM},\,\overrightarrow{IO})=\arccos \left(\dfrac{-1\cdot(-7)}{1\cdot\sqrt{634}} \right)=\arccos \left(\dfrac{7}{\sqrt{634}} \right)$.\\
			Khi đó, thời  gian mục tiêu di chuyển từ $M$ đến điểm $C$ là~\[{t_{mt}}=\dfrac{1\cdot\arccos \left(\dfrac{7}{\sqrt{634}} \right)}{\arccos \left( \dfrac{24}{7} \right) }\approx 1{,}0016 \text{ (giây)}.\]
			Thời gian viên  đạn bay đến $C$ là~${t_{vd}}=\dfrac{OC}{v}=\dfrac{\sqrt{634}-1}{3\sqrt{65}}\approx 0{,}99969$ giây.
		\end{itemchoice}
	} 
\end{ex}

% nguồn: Phan1_Lop_12_Hinh_hoc_Chuong_5
\begin{ex}%[ma:fa0e9cea]
	%[2H5V2-7] 
	 \immini{
	Cho hình chóp $SABCD$ có đáy $ABCD$ là~hình  vuông và $SA$ vuông góc với đáy, $SA=2a$,\break $AB=a$.
	\choiceTF 
		{\True  Thể tích khối chóp $SABCD$ bằng $\dfrac{2a^3}{3}$}
		{Khoảng cách giữa hai đường thẳng $SA$ và $BC$ bằng  $\dfrac{2a}{\sqrt{5}}$}
		{\True  $\overrightarrow{SA}+\overrightarrow{SC}=\overrightarrow{SB}+\overrightarrow{SD}$}
		{Số đo góc nhị diện $[B,SC,D]$ bằng $103,5^{\circ}$ (làm  tròn đến hàng phần chục)}
		} 
		{ 
		\begin{tikzpicture}[line join=round,line cap=round,>=stealth]
				\path 
					 (0,0) coordinate (A)
					 ($(A)+(-120:1.5)$) coordinate (B)
					 ($(A)+(0:4)$) coordinate (D)
					($(B)+(D)-(A)$) coordinate  (C)
					($(A)+(90:3)$) coordinate (S); 
				\draw 
					 (S)--(B)--(C)--(D)--cycle
					 (S)--(C);
				\draw[dashed] 
					 (B)--(A)--(D)
					 (S)--(A)--(C);
				\foreach \x/\g in  {A/135,B/180,C/-90,D/0,S/90}
				\fill (\x) circle  (1pt) node[shift={(\g:10pt)}] {$\x$};
			\end{tikzpicture}
		 }
	 \loigiai{
		\begin{center}
			\begin{tikzpicture}[line join=round,line cap=round,>=stealth]
					 \path
						 (0,0) coordinate (A)
						($(A)+(-120:1.5)$)  coordinate (B)
						($(A)+(0:4)$) coordinate  (D)
						($(B)+(D)-(A)$) coordinate (C) 
						($(A)+(90:3)$) coordinate (S) 
						($(A)!1/2!(C)$) coordinate (O) 
						($(S)!1/2!(C)$)  coordinate (M);
					\draw 
						(S)--(B)--(C)--(D)--cycle 
						 (S)--(C)
						(B)--(M)--(D); 
					\draw[dashed] 
						(B)--(A)--(D)--cycle 
						 (S)--(A)--(C);
					\foreach \x/\g in {A/135,B/180,C/-90,D/0,S/90,O/-90,M/45} 
					 \fill (\x) circle (1pt) node[shift={(\g:10pt)}] {$\x$};
				\end{tikzpicture}
		\end{center}
		\begin{itemchoice}
			\itemch 
			Diện tích đáy là~$S_{ABCD}=a^2$, thể tích khối  chóp là~$V_{SABCD}=\dfrac{1}{3}SA\cdot S_{ABCD}=\dfrac{1}{3}\cdot 2a\cdot a^2=\dfrac{2a^3}{3}$.
			\itemch 
			Ta có  $AB\perp SA,AB\perp BC\Rightarrow\mathrm{d}(SA,BC)=AB=a$.
			\itemch 
			Gọi $O$ là~giao điểm của $AC$ và $BD$ thì $O$ cũng là~trung điểm của $AC$ và $BD$. Khi đó  $\overrightarrow{SA}+\overrightarrow{SC}=\overrightarrow{SB}+\overrightarrow{SD}=2\overrightarrow{SO}$.
			 \itemch
			Ta  có $SB=\sqrt{SA^2+AB^2}=\sqrt{4a^2+a^2}=a\sqrt{5}$.\\
			Xét tam giác $SBC$ vuông tại  $B$ nên $SC=\sqrt{SB^2+BC^2}=\sqrt{5a^2+a^2}=a\sqrt{6}$.\\
			Trong mặt phẳng $(SBC)$  kẻ $BM\perp SC$, ta có $\heva{&BM\perp SC\\&BD\perp SC}\Rightarrow SC\perp(BDM)\Rightarrow SC\perp MD$, do vậy góc nhị diện $[B,SC,D]$ là~góc $\widehat{BMD}$.\\
			Xét tam giác $BDM$ có $MB=MD=\dfrac{BC\cdot BS}{SC}=\dfrac{a\sqrt{5}}{\sqrt{6}}$, $BD=a\sqrt{2}$ nên $$\cos\widehat{BMD}=\dfrac{BM^2+DM^2-BD^2}{2\cdot  BM\cdot DM}=-\dfrac{1}{5}.$$
			Do vậy  góc $\widehat{BMD}=101{,}5^{\circ}$.
		\end{itemchoice}
	} 
\end{ex}

% nguồn: Phan1_Lop_12_Hinh_hoc_Chuong_5
\begin{ex}%[ma:ee9c67e3]
	 %[2H5H2-7]
	 Trong không gian với hệ tọa~độ $Oxyz$, cho hai đường thẳng $\Delta_1\colon\heva{&x=-1+t_1 \\& y=0 \\& z=1-t_1}$, $\Delta_2\colon\heva{&x=1-t_2 \\& y=-t_2 \\ &z=-1.}$
	\choiceTF 
		{\True$\vec{u}_1=(1; 0; -1)$~là một vectơ chỉ phương của $\Delta_1$} 
		{$\vec{u}_2=(1; -1; 0)$~là một vectơ chỉ phương của $\Delta_2$} 
		{\True$\vec{u}_1 \cdot \vec{u}_2=-1$} 
		{Góc giữa  hai đường thẳng $\Delta_1\colon\heva{&x=-1+t_1 \\& y=0 \\& z=1-t_1}$, $\Delta_2\colon\heva{&x=1-t_2 \\& y=-t_2 \\ &z=-1}$~là $120^\circ$}
	\loigiai{ 
		\begin{itemchoice}
			\itemch  \textbf{Đúng.} Vì đường thẳng $\Delta_1\colon\heva{&x=-1+t_1 \\& y=0 \\& z=1-t_1}$ có một vectơ chỉ phương $\vec{u}_1=(1; 0; -1)$.
			\itemch \textbf{Sai.} Vì đường thẳng $\Delta_2\colon\heva{&x=1-t_2 \\& y=-t_2 \\ &z=-1}$ có một vectơ chỉ  phương $\vec{u}_3=(-1; -1; 0)$.
			\itemch  \textbf{Đúng.} Vì $\vec{u}_1 \cdot \vec{u}_3=1\cdot(-1)+0\cdot(-1)+(-1)\cdot0=-1$.
			\itemch \textbf{Sai.} Vì $\cos\left(\Delta_1,\Delta_2\right)=\dfrac{\left|  \vec{u}_1\cdot\vec{u}_3\right| }{\left|\vec{u}_1\right|\cdot\left|\vec{u}_3\right|}=\dfrac{\left| -1\right|}{\sqrt{1+0+1}\cdot\sqrt{1+1+0}}=\dfrac{1}{2}\Rightarrow \left(\Delta_1,\Delta_2\right) =60^\circ$.
		\end{itemchoice}
	 }
\end{ex}

% nguồn: Phan1_Lop_12_Hinh_hoc_Chuong_5
\begin{ex}%[ma:052a5e7e]
	Xét một hệ trục tọa độ $Oxyz$ được cho sẵn, đơn~vị trên mỗi trục là dm, mặt ngoài của một quả bóng được mô hình hóa bởi phương trình mặt cầu $\left( x-2 \right)^{2}+\left( y+1 \right)^{2}+\left( z+1 \right)^{2}=6$, quả bóng nằm yên trên sàn nhà. Người ta nhìn thấy một tấm ván ngã xuống đè lên quả bóng, phần giao của tấm ván và sàn nhà là đường thẳng $d$ có phương trình $\dfrac{x+2}{2}=\dfrac{y+1}{-3}=\dfrac{z}{1}$. Gọi $A$, $B$ lần lượt là hai tiếp điểm  của tấm ván, sàn nhà với quả bóng và $I$ là tâm quả bóng.
	Xét  tính đúng sai các mệnh đề sau: \\
	\choiceTF 
		{\True Quả bóng có tâm $I\left( 2;-1;-1\right) $ và bán kính $R=\sqrt{6}$}%[2H5N3-2] 
		{Khoảng cách từ tâm quả  bóng đến đường thẳng $d$ bằng $2\sqrt{6}$}%[2H5H2-6]
		{\True Nếu $\cos \widehat{AIB}$ bằng $\dfrac{a}{b}$ (phân số tối giản) thì giá  trị $a^2+b^2=82$}%[2H5V2-7]
		{Một con kiến bò từ vị trí $A$ đến vị trí $B$ trên  quả bóng với tốc độ $2~\mathrm{cm/s}$; thời gian ngắn nhất cho chuyến đi này là $21$ giây (làm tròn đến hàng đơn~vị)}%[2H5V3-4]
	\loigiai{ 
		\begin{center}
			\begin{tikzpicture}[scale=.75,font=\footnotesize]
					\def\r{2} 
					\path  (0:0) coordinate (I)
						(-150:2*\r) coordinate (K) 
						($(I)!.5!(K)$) coordinate  (O);
					\path[name path=ciro]  (I) circle (\r);
					\path[name path=cirom] (O)  let\p1=($(O)-(I)$) in circle ({veclen (\x1,\y1)});
					\path[name intersections={of  = ciro and cirom,by={G,B}}];
					\path 
						 ($(K)!1!-6:(G)$) coordinate (A)
						(intersection  of A--B and I--K) coordinate (H)
						(-3,-3) coordinate (S) 
						(0,2) coordinate  (X)
						($(X)!0.5!(B)$) coordinate (L); 
					\pgfmathanglebetweenpoints{\pgfpointanchor{L}{center}}{\pgfpointanchor{X}{center}} 
					 \let\StartAngle\pgfmathresult
					\draw (X) let \p1 = ($(L) - (B)$), 
					\n1 =  {veclen (\x1,\y1)} in arc (\StartAngle:{\StartAngle-180}:{\n1});
					\path ($(X)!0.5!(B)$) coordinate (Q); 
					 \pgfmathanglebetweenpoints{\pgfpointanchor{Q}{center}}{\pgfpointanchor{X}{center}}
					 \let\StartAngle\pgfmathresult
					\draw[dashed] (X) let \p1 = ($(Q) -  (B)$),
					\n1 = {veclen (\x1,\y1)} in  arc (\StartAngle:{\StartAngle+180}:{\n1});
					 \path
						($(S)!2.8!(K)$) coordinate (T) 
						($(B)!-1!(K)$)  coordinate (N)
						 ($(I)!(T)!(B)$) coordinate (Th)
						(barycentric cs:T=1,Th=-2)  coordinate (V)
						($(S)!0.5!(V)$) coordinate  (P)
						($(P)!-1!(T)$)  coordinate (U)
						($(K)!2!(A)$) coordinate  (E)
						($(S)!0.5!(E)$)  coordinate (J)
						($(J)!-1!(K)$)  coordinate (F)
						($(T)!0.5!(E)$)  coordinate (M)
						($(M)!-1!(K)$) coordinate (W) 
						($(K)!1.5!(T)$) coordinate  (C)
						($(K)!1.8!(S)$) coordinate  (D)
						(D) node[above right]{$d$} 
						(S)--(D)--([turn]90:7pt) coordinate  (xt)
						($(S)-(D)+(xt)$) coordinate  (yt)
						(intersection of  K--B and S--F) coordinate (R)
						 ($(T)!.73!(V)$) coordinate (Y);
					 \draw[dashed]
						(0:\r) arc (0:180:{\r} and  {\r/3})
						 (90:\r) arc (90:270:{\r/1.25} and {\r})
						(90:\r) arc (90:-270:{\r/1.25} and  {\r})
						 (I)--(O)
						(A)--(B) 
						 (A)--(I)--(B)
						(180:\r) arc (180:235:{\r} and  {\r/3})
						 (O)--(K)
						(K)--(R) 
						(T)--(Y); 
					\draw 
						(-1.2,-.55) arc  (-140:11:{\r/1.1} and {\r/3.06})
						(K)--(A) 
						(S)--(T) 
						 (V)--(U)--(S)
						(T)--(W)--(F)--(S) 
						 (T)--(C)
						 (S)--(D)
						 (Y)--(V)
						 (T)--(W)--(F)--(S)
						 (T)--(C)
						 (S)--(D)
						 (R)--(B);
					\foreach \x/\y/\z in  {T/W/F}{
						\path (\y)  pic[draw,angle radius=20pt,angle eccentricity=0.5,"$P$"]{angle = \x--\y--\z};
					} 
					\foreach \x/\y/\z in  {V/U/S}{
						\path (\y) pic[draw,angle radius=20pt,angle eccentricity=0.5,"$Q$"]{angle = \x--\y--\z}; 
					} 
					 \foreach \x/\y/\z in {I/H/B}{
						\path (\y) pic[draw,angle radius=5pt]{right  angle = \x--\y--\z};
					 }
					\foreach \x/\g  in {I/60,K/-150,B/-90,H/-120,A/20} \draw (\x) circle (1pt) +(\g:.3) node{$\x$};
				\end{tikzpicture}
		\end{center}
		\begin{enumerate}[a)]
			\item Mệnh đề  đúng.\\
			Mặt ngoài của  quả bóng là mặt cầu $(S)$ có tâm $I\left(2;-1;-1 \right) $ và bán kính $R=\sqrt{6}$.
			\item Mệnh đề  sai.\\
			Đường thẳng $d\colon \dfrac{x+2}{2}=\dfrac{y+1}{-3}=\dfrac{z}{1}$ qua $A\left(-2;-1;0 \right)$ và  có véc-tơ chỉ phương $\vec{u}_d=(2;-3;1)$.\\
			Ta có $\overrightarrow{AI}=(4;0;-1)$; $\left[\vec{u}_d{,}\overrightarrow{AI}\right]=\left( \begin{vmatrix}
			 -3 & 1 \\
			0 & -1 
		\end{vmatrix}; \begin{vmatrix}
		1 &  2 \\
		 -1 & 4
	\end{vmatrix}; \begin{vmatrix}
	2 &  -3 \\
	 4 & 0
	\end{vmatrix}\right) = \left(3;6;12 \right) $.\\
	Do đó  $\mathrm{\,d}(I,d)=\dfrac{\left|\left[\vec{u}_d {,} \overrightarrow{AI}\right] \right| }{\left|\vec{u}_d \right| }$\\
	$=\dfrac{\sqrt{3^2+6^2+12^2}}{\sqrt{2^2+(-3)^2+1^2}}=\dfrac{3\sqrt{6}}{2}$. 
	\item  Mệnh đề đúng.\\
	Gọi $K$ là hình chiếu của $I$ trên $d$ thì $KI=\dfrac{3\sqrt{6}}{2}$ và  $KA \perp IA$.\\
	Suy ra $\cos \widehat{AIK}=\dfrac{IA}{IK}=\dfrac{2}{3}$.\\ 
	Do vậy $\cos \widehat{AIB}=2\cos^2\widehat{AIK}-1 =  2\times \left(\dfrac{2}{3} \right)^2-1=-\dfrac{1}{9}=\dfrac{a}{b}$.\\
	 Suy ra $a^2+b^2=(-1)^2+9^2=82$.
	\item Mệnh  đề sai.\\
	Độ dài cung tròn bé nhất mà con kiến có thể đi là\\ $\ell_{\stackrel{\frown}{AB}}=R\times \widehat{AIB}=\sqrt{6}\times  \arccos\left( - \dfrac{1}{9}\right) \approx 4{,}12~\mathrm{dm}$.\\
	Thời gian tối thiểu để kiến đến nơi là $t=\dfrac{10\times \ell_{\stackrel{\frown}{AB}}}{2}\approx  21$ giây.
	\end{enumerate}
 }
\end{ex}

% nguồn: Phan1_Lop_12_Hinh_hoc_Chuong_5
\begin{ex}%[ma:08b62063]
	 %[2H5H2-8]
	Một máy bay đang di chuyển về phía sân bay. Tại thời điểm hiện tại, vị trí của máy bay là~$B(150;  150; 5000)$ (trong đó $5000$ m là~độ cao của máy bay so với mặt đất). Máy bay đang di chuyển thẳng tới sân bay với vận tốc $700$ km/h. Sân bay có tọa độ $C(0; 0; 0)$ và máy bay đang tiến dần đến vị trí hạ cánh tại sân bay.
	\choiceTF 
		{\True Phương trình tham số của đường thẳng mà máy bay di chuyển theo là~$\left\{\begin{aligned}&x=150-150 t \\ & y=150-150 t \\ & z=5000-5000 t
\end{aligned}\right.$}
{Khoảng cách từ vị  trí hiện tại của máy bay $B(150; 150; 5000)$ đến sân bay $C(0; 0; 0)$ là~$\sqrt{15250000} \approx 3905{,}6$ km}
{Với vận tốc của máy bay là~$700$ km/h,  thời gian để máy bay hạ cánh là~khoảng $5{,}5$ giờ}
{Nếu hệ thống kiểm soát không lưu yêu cầu liên lạc với máy bay khi nó còn cách  sân bay $40$ km thì khi máy bay ở vị trí $(6; 6; 200)$ nó còn cách sân bay là~$40$ km}
 \loigiai{
	\begin{itemchoice}
		\itemch Véc-tơ  chỉ phương của đường thẳng $BC$ là~$\overrightarrow{BC}=(-150;-150;-5000)$ do đó phương trình tham số là~$\heva{&x=150-150 t \\ &y=150-150 t \\& z=5000-5000 t.}$
		\itemch $BC=\sqrt{(150-0)^2+(150-0)^2+(5000-0)^2}=\sqrt{25045000} \approx 5004{,}5$. 
		\itemch  Thời gian hạ cánh là~$t=\dfrac{5004{,}5}{700} \approx 7{,}15$ giờ.
		\itemch Tọa độ máy bay tại thời điểm còn  cách sân bay $40$ km phải thỏa mãn điều kiện khoảng cách là~$40$ km.\\
		 Ta có khoảng cách từ $(6; 6; 200)$ đến $(0; 0; 0)$ là~$d=\sqrt{(6-0)^2+(6-0)^2+(200-0)^2}=\sqrt{40072} \approx 200{,}18$ km.
	\end{itemchoice}
} 
\end{ex}

% nguồn: Phan1_Lop_12_Hinh_hoc_Chuong_5
\begin{ex}%[ma:f38bdb5c]
	%[2H5H2-4] 
	Trong  không gian $Oxyz$, cho hai đường thẳng $\Delta_1\colon \dfrac{x+3}{2}=\dfrac{y+2}{1}=\dfrac{z}{-1}$ và $\Delta_2\colon \dfrac{x+3}{1}=\dfrac{y+5}{2}=\dfrac{z-1}{-1}$.
	 \choiceTF
		{\True Đường thẳng $\Delta_1$ đi qua điểm  $M_1(-3;-2;0)$ và có $\vec{u}_{1}=(2;1;-1)$~là véc-tơ chỉ phương}
		{\True Đường thẳng $\Delta_2$ đi  qua điểm $M_2(-3;-5;1)$ và có $\vec{u}_{2}=(1;2;-1)$~là véc-tơ chỉ phương}
		{\True Nếu $\vec{u}_{1}=(2;1;-1)$ và $\vec{u}_{2}=(1;2;-1)$ thì  $\left[\vec{u}_{1},\vec{u}_{2}\right]=(1;1;3)$}
		{Hai đường  thẳng $\Delta_1$ và $\Delta_2$ chéo nhau}
	 \loigiai{
	\begin{itemchoice}
		\itemch \textbf{Đúng}. Đường thẳng $\Delta_1$  đi qua điểm $M_1(-3;-2;0)$ và có $\vec{u}_{1}=(2;1;-1)$~là véc-tơ chỉ phương.
		\itemch \textbf{Đúng}. Đường thẳng $\Delta_2$ đi qua điểm $M_2(-3;-5;1)$ và có  $\vec{u}_{2}=(1;2;-1)$~là véc-tơ chỉ phương.
		\itemch \textbf{Đúng}. Nếu  $\vec{u}_{1}=(2;1;-1)$ và $\vec{u}_{2}=(1;2;-1)$ thì $\left[\vec{u}_{1},\vec{u}_{2}\right]=(1;1;3)$.
		\itemch \textbf{Sai}. Ta có $\overrightarrow{M_1M_2}=(0;-3;1)$ nên $\left[\vec{u}_{1},\vec{u}_{2}\right]\cdot\overrightarrow{M_1M_2}=0$.  \\
		Vậy đường thẳng  $\Delta_1$ và $\Delta_2$ đồng phẳng.
	\end{itemchoice}
	} 
\end{ex}

% nguồn: Phan1_Lop_12_Hinh_hoc_Chuong_5
\begin{ex}%[ma:ca1044e7]
	 %[2H5H3-3]
	Trong không gian $Oxyz$  (đơn~vị trên mỗi trục tính theo mét), một ngọn hải đăng được đặt ở vị trí $I(17;20;45)$. Biết rằng ngọn hải đăng đó được thiết kế với bán kính phủ sáng~là $4\rm\,km$.
	 \choiceTF
		{\True Phương trình mặt cầu để mô tả  ranh giới bên ngoài của vùng phủ sáng trên biển của hải đăng~là $(x-17)^2+(y-20)^2+(z-45)^2=4000^2$}
		{Nếu người đi biển ở vị  trí $M(18;21;50)$ thì không thể nhìn thấy được ánh sáng từ ngọn hải đăng}
		{Nếu người đi biển ở vị trí $N(4019;21;44)$ thì có thể nhìn thấy được  ánh sáng từ ngọn hải đăng}
		{\True Nếu hai người đi biển ở vị trí đó có thể nhìn thấy được ánh sáng từ ngọn hải đăng thì khoảng cách giữa hai người đó không  quá $8\rm\,km$}
	 \loigiai{
		\begin{itemchoice}
			\itemch \textbf{Đúng}.\\ 
			Mặt cầu $(S)$ tâm $I(17;20;45)$  bán kính $R=4000\rm\,m$\\
			 $(S)\colon (x-17)^2+(y-20)^2+(z-45)^2=4000^2$.
			\itemch  \textbf{Sai}.\\
			Ta có $\heva{&I(17;20;45)\\& M(18;21;50)} \Rightarrow  \overrightarrow{IM}=(1;1;5) \Rightarrow IM=3\sqrt{3}$.\\
			Vì $3\sqrt{3}<m<4\rm\,km$ nên trong vùng phủ  sáng.
			 \itemch \textbf{Sai}.\\
			 Ta có $\heva{&I(17;20;45)\\ &N(4019;21;44)} \Rightarrow \overrightarrow{IN}=(4002;1;-1) \Rightarrow IN\approx4002$.\\
			Ngoài vùng phủ sáng. 
			\itemch  \textbf{Đúng}.\\
			Vì  bán kính phủ sáng $R=4\rm\,km$ nên có đường kính $d=8\rm\,km$.\\
			Do đó  khoảng cách hai người này không quá $d=8\rm\,km$.
		\end{itemchoice}
	} 
\end{ex}

% nguồn: Phan1_Lop_12_Hinh_hoc_Chuong_5
\begin{ex}%[ma:633c8fd4]
	%[2H5N1-2]%[2H5H1-6] 
	Trong không gian $Oxyz$,  cho hai mặt phẳng $(P)\colon y=0$,\break
	$(Q)\colon \sqrt{3}x-y-2024=0$. Xét các vectơ $\vec{n}_1=(0;1;0)$,  $\vec{n}_2=\big(\sqrt{3};-1;0\big)$.
	 \choiceTF
		{\True $\vec{n}_1$~là một vectơ pháp tuyến  của mặt phẳng $(P)$}
		{$\vec{n}_2$ không~là vectơ pháp tuyến của mặt phẳng $(Q)$} 
		{\True  $\vec{n}_1\cdot\vec{n}_2=-1$}
		{Góc  giữa hai mặt phẳng $(P)$ và $(Q)$ bằng $30^\circ$}
	 \loigiai{
		\begin{itemchoice}
			\itemch \textbf{Đúng}. $\vec{n}_1$~là một  vectơ pháp tuyến của mặt phẳng $(P)$.\\
			Ta có  $(P)\colon y=0 \Leftrightarrow 0x+1y+0z=0$.\\
			 vectơ pháp tuyến $\vec{n}_1=(0;1;0)$.
			\itemch \textbf{Sai}. $\vec{n}_2$ không~là vectơ pháp tuyến  của mặt phẳng $(Q)$.\\
			Ta có $(Q)\colon \sqrt{3}x-y-2024=0 \Leftrightarrow  \sqrt{3}x-y+0z-2024=0$.\\
			 vectơ pháp tuyến $\vec{n}_2=\left(\sqrt{3};-1;0\right)$.
			\itemch \textbf{Đúng}. $\vec{n}_1\cdot\vec{n}_2=-1$.\\ 
			Ta  có $\heva{\vec{n}_1&=(0;1;0)\\ \vec{n}_2&=\left(\sqrt{3};-1;0\right)} \Rightarrow \vec{n}_1\cdot\vec{n}_2=0\cdot\sqrt{3}+1\cdot(-1)+0\cdot0=-1$.
			\itemch \textbf{Sai}.  Góc giữa hai mặt phẳng $(P)$ và $(Q)$ bằng $30^\circ$.\\
			Gọi $\varphi$~là góc giữa  hai mặt phẳng $(P)$ và $(Q)$
			\begin{align*}
				 \cos\varphi=\left|\cos\left(\vec{n}_1,\vec{n}_2\right)\right|
				 &=\dfrac{\left|\vec{n}_1\cdot\vec{n}_2\right|}{|\vec{n}_1|\cdot|\vec{n}_2|}\\
				& =\dfrac{\left|0\cdot\sqrt{3}+1\cdot(-1)+0\cdot0\right|}{\sqrt{0^2+1^2+0^2}\cdot\sqrt{\left(\sqrt{3}\right)^2+(-1)^2+0^2}}\\ &=\dfrac{1}{2}\cdot 
			\end{align*}
			Suy ra  góc $\varphi=60^\circ$.
		\end{itemchoice}
	} 
\end{ex}
